\subsection{特征标与共轭类}

令 \(\varrho\) 为 G 在有限维希尔伯特空间 E 中的酉表示。\(\varrho\) 的特征标满足 \(|\chi_{\varrho}| \leqslant \dim(E)\) 。令 \((e_i)_{i \in I}\) 为 E 的标准正交基；特征标 \(\chi_{\varrho}\) 是对角矩阵系数 \(g \mapsto \langle e_i | \varrho(g)e_i \rangle\)（属于 \(\varrho\)）之和。特别地，\(\varrho\) 的特征标是 G-有限函数；若 \(\varrho\) 不可约，则 \(\chi_{\varrho} \in \varrho_{\mathrm{G}}(\varrho)\) 。

有如下公式

\[
\chi_ {\varrho_ {1} \oplus \varrho_ {2}} = \chi_ {\varrho_ {1}} + \chi_ {\varrho_ {2}}, \quad \chi_ {\varrho_ {1} \otimes \varrho_ {2}} = \chi_ {\varrho_ {1}} \chi_ {\varrho_ {2}}, \quad \chi_ {\breve {\varrho}} = \chi_ {\overline {{\varrho}}} = \overline {{\chi}} _ {\varrho}
\]

（参见 A, VIII, 第 388--389 页）。

命题 8. --- 有

\begin{enumerate}
\def\labelenumi{(\arabic{enumi})}
\tightlist
\item
  \(\chi_{\pi}*\chi_{\sigma} = 0\) 对于 \(\pi, \sigma\) 均属于 \(\widehat{\mathrm{G}}\) 且 \(\pi \neq \sigma\) 的情形成立，
\end{enumerate}

\[
\chi_ {\pi} * \chi_ {\pi} = \frac {1}{\dim (\pi)} \chi_ {\pi} \quad\text{对于每个 }\pi\in\widehat {\mathrm{G}}.\tag{2}
\]

令 \(\pi\) 和 \(\sigma\) 为 G 的不可约表示。有

\[
\dim (\pi) \left(\chi_ {\pi} * \chi_ {\sigma}\right) = \dim (\pi) \boldsymbol {\gamma} _ {G} \left(\chi_ {\pi}\right) \chi_ {\sigma}
\]

（V, 第 407 页引理 4）。函数 \(\dim(\pi)\gamma_{\mathrm{G}}(\chi_{\pi})\chi_{\sigma}\) 是将 \(\chi_{\sigma}\) 投影到 \(\overline{\pi}\)-同型分量 \(\gamma_{\mathrm{G}}\) 上的正交投影（V, 第 465 页推论 2），即投影到 \(\varrho_{\mathrm{G}}(\pi)\) 上（V, 第 463 页推论 4）。由于 \(\chi_{\sigma}\) 属于 \(\varrho_{\mathrm{G}}(\sigma)\)，结论由 V, 第 462 页定理 1 得出。

引理 2. --- 等价关系“\(x \in G\) 且 \(y \in G\) 且 \(x\) 与 \(y\) 共轭”在 G 中的图像是闭集。

事实上，该图像是从紧空间 G × G 到自身、由 \((x, y) \mapsto (x, yxy^{-1})\) 定义的连续映射的像。

记 G \(^{\sharp}\) 为赋予商拓扑的 G 的共轭类空间；由引理 2 以及 TG, I, 第 78 页命题 8，它是紧空间。令 \(\varpi\) : G → G \(^{\sharp}\) 为典范投影。由从 \(\mathcal{C}(G^{\sharp})\) 到 \(\mathcal{C}(G)\) 的映射 \(f \mapsto f \circ \varpi\)，将星代数 \(\mathcal{C}(G^{\sharp})\) 识别为 \(\mathcal{C}(G)\) 中由连续中心函数组成的星子代数。

G \(^{\sharp}\) 上的测度与 G 上的中心测度相对应（V, 第 402 页定义 1）。

于是，在 \(\mathcal{C}(\mathrm{G}^{\sharp})\) 上由 \(f\mapsto\int_{\mathrm{G}}f\) 定义的线性泛函是 \(\mathrm{G}^{\sharp}\) 上质量为 1 的正测度，记为 \(\mu^{\sharp}\)。对于每个 \(p\in[1,+\infty]\)，记 \(\mathcal{L}^{p}(\mathrm{G}^{\sharp})\)（分别记 \(\mathrm{L}^{p}(\mathrm{G}^{\sharp})\)）为 \(\mathcal{L}_{\mathbf{C}}^{p}(\mathrm{G}^{\sharp},\mu^{\sharp})\)（分别为 \(\mathrm{L}_{\mathbf{C}}^{p}(\mathrm{G}^{\sharp},\mu^{\sharp})\)）。我们将 \(\mathrm{L}^{p}(\mathrm{G}^{\sharp})\) 识别为 \(\mathcal{C}(\mathrm{G}^{\sharp})\) 在 \(\mathrm{L}^{p}(\mathrm{G})\) 中的闭包；特别地，它是 \(\mathrm{L}^{p}(\mathrm{G})\) 的闭子空间。

我们还记 \(\Theta(G^{\sharp}) = \Theta(G) \cap \mathcal{C}(G^{\sharp})\) 为 G 的中心矩阵系数空间。它是 \(\mathcal{C}(G^{\sharp})\) 的含单位对合子代数。

当 G 为李群时，在 LIE, IX, 第 71 页中，空间 \(\Theta(G^{\sharp})\) 也记为 \(Z\Theta(G)\)。

引理 3. --- 令 \(\pi\) 为 G 的不可约酉表示。向量空间 \(\mathrm{L}^2 (\mathrm{G}^\sharp)\cap \varrho_{\mathrm{G}}(\pi)\) 是一维的，并由 \(\pi\) 的特征标生成。

考虑 G 的酉表示 \(\sigma\)，它作用在空间 \(\varrho_{\mathrm{G}}(\pi)\) 上，并由 \(\sigma(g) = \varrho_{\mathrm{G}}(g, g)\) 定义。由于 G × G 的酉表示 \(\varrho_{\mathrm{G}}(\pi)\) 同构于 \(\overline{\pi} \boxtimes \pi\)（V, 第 462 页命题 5），其特征标为

\[
\chi_ {\sigma} (g) = \chi_ {\overline {{\pi}} \boxtimes \pi} (g, g) = | \chi_ {\pi} (g) | ^ {2}
\]

对所有 \(g \in G\) 都成立。\(\mathrm{L}^2 (\mathrm{G}^\sharp) \cap \varrho_{\mathrm{G}}(\pi)\) 是该表示的不变元素子空间，且其维数等于

\[
\int_ {\mathrm{G}} \chi_ {\sigma} (g) d \mu (g) = \int_ {\mathrm{G}} | \chi_ {\pi} (g) | ^ {2} d \mu (g) = 1
\]

（V, 第 466 页推论 3 以及第 459 页推论 3）。由于 \(\pi\) 的特征标是 \(L^2(G^\sharp) \cap \varrho_G(\pi)\) 的非零元素，它就是该空间的一个基。

命题 9. --- G 的不可约酉表示的特征标族 \((\chi_{\pi})_{\pi \in \widehat{\mathbf{G}}}\) 是 \(\mathrm{L}^2 (\mathrm{G}^\sharp)\) 的标准正交基。

根据彼得–外尔定理（V, 第 462 页定理 1），不变元素构成闭子空间 \(L^{2}(G^{\sharp})\)，它是 \(L^{2}(G)\) 的子空间，由 G 的酉表示 \(g \mapsto \varrho_{\mathrm{G}}(g, g)\) 的不变元素构成；它是各子空间 \(L^{2}(G^{\sharp}) \cap \varrho_{\mathrm{G}}(\pi)\)（\(\pi \in \widehat{G}\)）的希尔伯特直和。因此命题由前一引理和 V, 第 459 页推论 3 得出。

推论 1. --- 向量空间 \(\Theta(\mathrm{G}^{\sharp})\) 在 \(\mathcal{C}(\mathrm{G}^{\sharp})\) 和 \(\mathrm{L}^{2}(\mathrm{G}^{\sharp})\) 中都是稠密的。

第二个断言由命题 9 得出。对于第一个断言，考虑 G 的线性表示 \(\varrho\)，它作用在巴拿赫空间 \(\mathcal{C}(\mathrm{G})\) 上，并由下式定义：

\[
\varrho (g) f (x) = f (g ^ {- 1} x g)
\]

对所有 \(f \in \mathcal{C}(\mathrm{G})\) 及每个 \(x \in \mathrm{G}\) 都成立。它是连续等距表示，而 \(\mathcal{C}(\mathrm{G}^{\sharp})\) 是该表示的不变元素空间。因此，连续投影 \(p = \varrho(1)\) 作用在 \(\mathcal{C}(\mathrm{G})\) 上，且其像为 \(\mathcal{C}(\mathrm{G}^{\sharp})\)（V, 第 466 页推论 3）；其范数满足 \(\leqslant 1\)。

令 \(f \in \mathcal{C}(\mathrm{G}^{\sharp})\) 且 \(\varepsilon > 0\)。存在 \(\widetilde{f} \in \Theta(\mathrm{G})\)，使得 \(\|f - \widetilde{f}\| \leqslant \varepsilon\)（V, 第 462 页命题 4），于是 \(\|f - p(\widetilde{f})\| = \|p(f - \widetilde{f})\| \leqslant \varepsilon\)；由于 \(p(\widetilde{f}) \in \mathcal{C}(\mathrm{G}^{\sharp})\)，我们得到 \(\Theta(\mathrm{G}^{\sharp})\) 在 \(\mathcal{C}(\mathrm{G}^{\sharp})\) 中稠密。

记 R(G) 为 G 在分离的有限维复向量空间中的连续表示所成的格罗滕迪克环（参见 LIE, IX, 第 70 页和 A, VIII, 第 182 页；将其应用于 G 在分离的有限维复向量空间中的连续表示这一加性类，并将其视为 \(\mathbf{C}[G]\)-模）。由于 G 在有限维分离拓扑向量空间中的每个连续线性表示都是半单的（V, 第 457 页命题 1），阿贝尔群 R(G) 是自由的，而不可约酉表示 \(\pi \in \widehat{\mathbf{G}}\) 的类构成一个基（A, VIII, 第 186 页命题 7）。

推论 2. --- a) G 的不可约表示的特征标族是 \(\Theta(G^{\sharp})\) 的一个基；

\begin{enumerate}
\def\labelenumi{\alph{enumi})}
\setcounter{enumi}{1}
\tightlist
\item
  映射 \(u: \mathrm{R}(\mathrm{G}) \otimes_{\mathbf{Z}} \mathbf{C} \to \Theta(\mathrm{G}^{\sharp})\)，满足 \(u(\pi \otimes 1) = \chi_{\pi}\)（对每个 \(\pi \in \widehat{\mathrm{G}}\)），是 \(\mathbf{C}\) 上代数的同构。
\end{enumerate}

第一个断言由命题 9 和推论 1 得出。

由于表示 \(\pi \in \widehat{\mathbf{G}}\) 的类构成自由 \(\mathbf{Z}\)-模 R(G) 的一个基，映射 \(u\) 定义良好。它是 \(\mathbf{C}\)-代数的态射，并由 \(a\) 可知它是同构。

推论 3. --- 令 H 为局部紧拓扑群，且 G 是 H 的紧子群。令 \(\varrho\) 为 H 在希尔伯特空间 E 中的酉表示。若 \(\pi\)-同型分量（在 \(\varrho\) 限制到 G 后）对所有 \(\pi \in \widehat{\mathbf{G}}\) 都是有限维的，则 H 的表示 \(\varrho\) 是半单的，并且 H 的每个不可约酉表示在 \(\varrho\) 中的重数都是有限的。

记 \(\sigma\) 为 \(\varrho\) 限制到 H 的紧子群 G 后的表示。令 \(\pi\) 为 G 的不可约表示。自同态 \(\sigma(\chi_{\pi}) \in \mathcal{L}(\mathrm{E})\) 是有限秩的，因为其像是 \(\overline{\pi}\)-同型分量，即 \(\sigma\) 中的相应子空间（V, 第 465 页推论 2），按假设它是有限维的。因此，自同态 \(\sigma(f)\) 对每个 \(f \in \Theta(\mathrm{G}^{\sharp})\) 都是有限秩的；对每个 \(f \in \mathcal{C}(\mathrm{G}^{\sharp})\)，它都是紧的（V, 第 468 页推论 1 以及 III, 第 4 页命题 2 的推论）。

记 \(j\) 为从 G 到 H 的典范嵌入。这是一个连续映射，并且对于 G 上的每个测度 \(\nu\) 都是 \(\nu\)-真映射（INT, V, 第 69 页，第 6 节第 1 号注 1）。令 \(\mathfrak{B}\) 为 G 中元素 \(e\) 的紧邻域滤子。对于每个 V \(\in\) \(\mathfrak{B}\)，G 上存在一个连续正中心函数 \(f_{\mathrm{V}}\)，其支撑包含于 V 且在 G 上的积分等于 1。令 \(\beta_{\mathrm{V}}\) 为像测度 \(j(f_{\mathrm{V}} \cdot \mu)\)；它是 H 上具有紧支撑的正测度，并满足 \(\varrho(\beta_{\mathrm{V}}) = \sigma(f_{\mathrm{V}})\)（INT, V, 第 71 页，第 6 节第 2 号定理 1）。根据上述结果，在 \(\mathcal{M}_{c}(\mathrm{H})\) 上，由 \(\mathfrak{B}\) 经映射 \(\mathrm{V} \mapsto \beta_{\mathrm{V}}\) 得到的滤子基满足 V, 第 402 页命题 2 的条件，断言得证。

推论 4（外尔等分布判据）

令 M 为 G 上的中心正测度集，满足 \(\nu(G)=1\) 对所有 \(\nu\in M\) 都成立。令 \(\mathfrak{F}\) 为 M 上的滤子。滤子 \(\mathfrak{F}\) 弱收敛到测度 \(\mu^{\sharp}\)，其中测度定义在 \(G^{\sharp}\) 上，其充要条件是：对每个非平凡不可约酉表示 \(\pi\)，都有

\[
\lim _ {\nu , \mathfrak {F}} \int_ {\mathrm{G}} \chi_ {\pi} (x) d \nu (x) = 0.
\]

由于 \(\nu (\mathrm{G}^{\sharp}) = 1 = \mu^{\sharp}(\mathrm{G}^{\sharp})\) 对 \(\nu \in \mathrm{M}\) 成立，上述假设意味着

\[
\lim _ {\nu , \mathfrak {F}} \int_ {\mathrm{G} ^ {\sharp}} \chi_ {\pi} (x) d \nu (x) = \int_ {\mathrm{G} ^ {\sharp}} \chi_ {\pi} (x) d \mu^ {\sharp} (x)
\]

对每个表示 \(\pi\in\widehat{\mathrm{G}}\) 都成立，因而它等价于如下条件

\[
\lim _ {\nu , \mathfrak {F}} \int_ {\mathrm{G} ^ {\sharp}} f (x) d \nu (x) = \int_ {\mathrm{G} ^ {\sharp}} f (x) d \mu^ {\sharp} (x)
\]

由于线性性，对每个函数 \(f \in \Theta(\mathrm{G}^{\sharp})\) 都成立。由于空间 \(\Theta(\mathrm{G}^{\sharp})\) 在 \(\mathcal{C}(\mathrm{G}^{\sharp})\) 中稠密（推论 1），所以该假设等价于滤子 \(\mathfrak{F}\) 收敛到 \(\mu^{\sharp}\)，并且在 \(\mathcal{M}(\mathrm{G}^{\sharp})\) 中收敛；该空间赋予由 \(\mathcal{C}(\mathrm{G})\) 上逐点收敛定义的拓扑。由于 G 紧，这一拓扑与弱收敛拓扑一致（参见 INT, III, 第 59 页，第 1 节第 9 号）。

\subsection{傅里叶余变换}

在集合 \(\widehat{\mathbf{G}}\) 上赋予离散拓扑。记 \(\mathrm{F}(\widehat{\mathbf{G}})\) 为各 \(\operatorname{End}(\mathrm{E}_{\pi})\) 的乘积代数，其中 \(\pi\) 遍历 \(\widehat{\mathbf{G}}\)，记 \(\mathrm{F}_b(\widehat{\mathbf{G}})\) 为各 \(\operatorname{End}(\mathrm{E}_{\pi})\) 的乘积星代数（I, 第 103 页例 5）；这就是由族 \((u_{\pi})_{\pi \in \widehat{\mathbf{G}}}\) 组成且满足 \(\sup \| u_{\pi} \| < +\infty\) 的集合。

\[
\pi \in \widehat {\mathbf {G}}
\]

记 \(\mathrm{F}_{0}(\widehat{\mathrm{G}})\) 为 \(\mathrm{F}_{b}(\widehat{\mathrm{G}})\) 的闭星子代数，它由族 \((u_{\pi})_{\pi\in\widehat{\mathrm{G}}}\) 组成，这些族满足 \(\|u_{\pi}\|\) 在无穷远处趋于 0。

令 \(\nu \in \mathcal{M}^1 (\mathrm{G})\)。对于每个 \(\pi \in \widehat{\mathrm{G}}\)，有 \(\| \pi (\nu)\| \leqslant \| \nu \|\)，所以族 \((\pi (\nu))_{\pi \in \widehat{\mathrm{G}}}\) 属于 \(\mathrm{F}_b(\widehat{\mathrm{G}})\)。

定义 1. --- 对每个测度 \(\nu \in \mathcal{M}^{1}(\mathrm{G})\)，将族 \((\pi (\nu))_{\pi \in \widehat{\mathrm{G}}}\) 视为 \(\mathrm{F}_b(\widehat{\mathrm{G}})\) 中的元素，记为 \(\overline{\mathcal{F}}_{\mathrm{G}}(\nu)\)。这样定义的从 \(\mathcal{M}^1 (\mathrm{G})\) 到 \(\mathrm{F}_b(\widehat{\mathrm{G}})\) 的映射称为 G 的傅里叶余变换，而 \(\overline{\mathcal{F}}_{\mathrm{G}}(\nu)\) 称为测度 \(\nu\) 的傅里叶余变换。

对于 \(f \in \mathrm{L}^1(\mathrm{G})\)，记 \(\overline{\mathcal{F}}_{\mathrm{G}}(f) = \overline{\mathcal{F}}_{\mathrm{G}}(f \cdot \mu)\)。

对于任意表示 \(\pi \in \widehat{\mathrm{G}}\)，映射 \(\nu \mapsto \pi(\nu)\) 是从 \(\mathcal{M}^{1}(\mathrm{G})\) 到 \(\operatorname{End}(\mathrm{E}_{\pi})\) 的含单位对合巴拿赫代数态射（V, 第 401 页引理 1）。因此，傅里叶余变换是从 \(\mathcal{M}^{1}(\mathrm{G})\) 到 \(\mathrm{F}_b(\widehat{\mathrm{G}})\) 的含单位对合巴拿赫代数态射。

\subsection{傅里叶变换}

沿用前一号的记号。

令 \(\pi \in \widehat{\mathbf{G}}\)。向量空间 \(\operatorname{End}(\mathrm{E}_{\pi})\) 赋予希尔伯特空间结构，其内积为

\[
\langle u _ {1} \mid u _ {2} \rangle = \dim (\pi) \operatorname{Tr} (u _ {1} ^ {*} u _ {2}) = \dim (\pi) \operatorname{Tr} (u _ {2} u _ {1} ^ {*})
\]

其中 \(u_{1}\)、\(u_{2}\) 属于 \(\operatorname{End}(\mathrm{E}_{\pi})\)（参见 EVT, V, 第 52 页定理 1）。

记 \(\|u\|_{2} = \sqrt{\langle u|u\rangle}\) 为 \(\operatorname{End}(\operatorname{E}_{\pi})\) 中元素 u 的范数，其中 \(\pi \in \widehat{G}\)。对于每个 \(g \in G\)，有 \(\|\pi(g)\|_{2} = \dim(\pi)\)，因为 \(\pi(g)\) 是酉的。

此处记作 \(\|u\|_{2}\) 的范数，与 EVT, V, 第 52 页定义的范数相差因子 \(\dim(\pi)\)；后者定义在 \(E_{\pi}\) 上的希尔伯特–施密特算子空间中。

引理 4. --- 令 \(\pi\) 为 G 的不可约表示。映射 \(f \mapsto \pi(f)\) 通过取子空间定义了从 \(\varrho_{\mathrm{G}}(\overline{\pi})\) 到 \(\operatorname{End}(\operatorname{E}_{\pi})\) 的等距同构。

置 \(\varepsilon_{\pi}(g,h)u = \pi(g)\circ u\circ\pi(h^{-1})\)，其中 \((g,h)\in\mathrm{G}\times\mathrm{G}\) 且 \(u\in\mathrm{End}(\mathrm{E}_{\pi})\)。映射 \(\varepsilon_{\pi}\) 是 \(\mathrm{G}\times\mathrm{G}\) 在 \(\mathrm{End}(\mathrm{E}_{\pi})\) 中的连续表示。它是酉的。事实上，由于 \(\pi\) 本身是酉的，有

\[
\begin{array}{c} \| \varepsilon_ {\pi} (g, h) u \| _ {2} ^ {2} = \dim (\pi) \operatorname{Tr} \Big ((\pi (g) u \pi (h ^ {- 1})) ^ {*} \pi (g) u \pi (h ^ {- 1}) \Big) \\ = \dim (\pi) \operatorname{Tr} (\pi (h) u ^ {*} u \pi (h) ^ {- 1}) = \dim (\pi) \operatorname{Tr} (u ^ {*} u) = \| u \| _ {2} ^ {2} \end{array}
\]

对所有 \((g,h)\in \mathrm{G}\times \mathrm{G}\) 及每个 \(u\in \operatorname {End}(\operatorname {E}_{\pi})\) 都成立。

映射 \(\Psi\)，定义为 \(f\mapsto\pi(f)\)，是一个 \((G\times G)\)-态射，从 \(\varrho_{G}(\overline{\pi})\) 到 \(\operatorname{End}(E_{\pi})\)，因为

\[
\pi (\varrho_ {\mathrm{G}} (g, h) f) = \int_ {\mathrm{G}} f (g ^ {- 1} x h) \pi (x) d \mu (x) = \pi (g) \pi (f) \pi (h ^ {- 1})
\]

对所有 \((g,h)\in\mathrm{G}\times\mathrm{G}\) 及每个 \(f\in\varrho_{\mathrm{G}}(\overline{\pi})\) 都成立。由于 \(\varrho_{\mathrm{G}}(\overline{\pi})\) 是不可约表示（V, 第 462 页定理 1，\(a\)），存在 \(\lambda\in\mathbf{C}^{*}\)，使得映射 \(\lambda\Psi\) 为零映射或等距映射（V, 第 388 页推论 5，\(a\)）。令 \(f=\dim(\pi)\overline{\chi}_{\pi}\in\varrho_{\mathrm{G}}(\overline{\pi})\)。得到 \(\pi(f)=1_{\mathrm{E}_{\pi}}\)（V, 第 465 页推论 2），从而 \(\|\pi(f)\|_{2}=\dim(\pi)=\|f\|\)（V, 第 459 页推论 3）。因此，映射 \(\Psi\) 是等距的。由于 \(\varrho_{\mathrm{G}}(\overline{\pi})\) 与 \(\operatorname{End}(\mathrm{E}_{\pi})\) 维数相同，\(\Psi\) 是等距同构。

记 \(\mathrm{F}^{2}(\widehat{\mathrm{G}})\) 为希尔伯特空间 \(\mathrm{End}(\mathrm{E}_{\pi})\) 的希尔伯特直和。元素 \(x \in \mathrm{F}^{2}(\widehat{\mathrm{G}})\) 的范数仍记为 \(\|x\|_{2}\)。

在 LIE, IX, 第 79 页中，该空间记为 \(\mathrm{L}^{2}(\widehat{\mathrm{G}})\)；这里我们不采用这一记号，以免与空间 \(\ell^{2}(\widehat{\mathrm{G}})\) 混淆。令 \((u_{\pi})_{\pi\in\widehat{\mathrm{G}}}\) 为 \(\mathrm{F}^{2}(\widehat{\mathrm{G}})\) 的元素。由于

\[
\sum_ {\pi \in \widehat {G}} \| u _ {\pi} \| _ {2} ^ {2} = \sum_ {\pi \in \widehat {G}} \dim (\pi) \operatorname{Tr} (u _ {\pi} ^ {*} u _ {\pi}) <   + \infty
\]

并且 \(\|u_{\pi}\|^{2}\leqslant\mathrm{Tr}(u_{\pi}^{*}u_{\pi})\)（参见 EVT, V, 第 52 页公式 (33)），在 \(\mathcal{L}(\mathrm{E}_{\pi})\) 中自同态 \(u_{\pi}\) 的范数在无穷远处趋于 0。因此可以将 \(\mathrm{F}^{2}(\widehat{\mathrm{G}})\) 识别为 \(\mathrm{F}_{0}(\widehat{\mathrm{G}})\) 的子空间。

令 \(\pi \in \widehat{\mathbf{G}}\) 且 \(u \in \mathrm{End}(\mathrm{E}_{\pi})\)。记 \(\mathcal{F}_{\pi}(u)\) 为 G 上由 \(\mathcal{F}_{\pi}(u)(g) = \langle \pi(g) | u \rangle = \dim(\pi) \mathrm{Tr}(\pi(g)^{*} u)\) 定义的函数（对所有 \(g \in \mathbf{G}\)）。它是 G 上的连续函数。若 G 是紧实李群，则函数 \(\mathcal{F}_{\pi}(u)\) 在 G 上解析（LIE, III, 第 8 节第 1 号定理 1）。

由于 G 是紧的，可以将 \(L^{2}(G)\) 识别为 \(L^{1}(G)\) 的子空间。

定理 2. --- G 的傅里叶余变换通过取子空间诱导出从 L²(G) 到 F²(Ĝ) 的等距同构。其逆 \(\mathcal{F}_{\mathrm{G}}\) 将 F²(Ĝ) 中的元素 \((u_{\pi})_{\pi \in \widehat{\mathrm{G}}}\) 对应于级数之和

\[
\sum_ {\pi \in \widehat {\mathrm{G}}} \mathcal {F} _ {\pi} (u _ {\pi}),
\]

该级数在 \(\mathrm{L}^2 (\mathrm{G})\) 中收敛。

根据彼得–外尔定理（V, 第 462 页定理 1），希尔伯特空间 \(L^{2}(G)\) 是各空间 \(\varrho_{\mathrm{G}}(\overline{\pi})\)（\(\pi \in \widehat{G}\)）的希尔伯特直和。对于每个 \(\pi \in \widehat{G}\)，线性映射 \(f \mapsto \pi(f)\) 从 \(\varrho_{\mathrm{G}}(\overline{\pi})\) 到 \(\operatorname{End}(\operatorname{E}_{\pi})\) 是等距同构（引理 4）。因此，傅里叶余变换在 \(L^{2}(G)\) 上的限制定义了从 \(L^{2}(G)\) 到 \(F^{2}(\widehat{G})\) 的等距同构。

令 \(f \in \mathrm{L}^2(\mathrm{G})\)。令 \(\pi \in \widehat{\mathrm{G}}\)，并令 \(f_{\overline{\pi}} \in \mathcal{C}(\mathrm{G})\) 为 \(f\) 在 \(\varrho_{\mathrm{G}}(\overline{\pi})\) 上的正交投影（V, 第 462 页定理 1）。由于该空间是 \(\overline{\pi}\)-同型分量，来自 \(\pmb{\delta}_{\mathrm{G}}\)（V, 第 463 页推论 4），由 V, 第 465 页推论 2，有 \(f_{\overline{\pi}} = \dim(\pi)\pmb{\delta}_{\mathrm{G}}(\chi_{\pi})(f)\)。于是对每个 \(x \in \mathrm{G}\)，有

\[
\begin{array}{r l} & {f _ {\overline {{\pi}}} (x) = \dim (\pi) \int_ {\mathrm{G}} \chi_ {\pi} (g) (\pmb {\delta} _ {\mathrm{G}} (g) f) (x) d \mu (g)} \\ & {\qquad = \dim (\pi) \int_ {\mathrm{G}} \chi_ {\pi} (g) f (x g) d \mu (g)} \\ & {\qquad = \dim (\pi) \int_ {\mathrm{G}} \chi_ {\pi} (x ^ {- 1} y) f (y) d \mu (y)} \\ & {\qquad = \dim (\pi) \mathrm{Tr} (\pi (x ^ {- 1}) \pi (f)) = \langle \pi (x) | \pi (f) \rangle ,} \end{array}
\]

即 \(f_{\overline{\pi}} = \mathcal{F}_{\pi}(\pi(f))\)。

由于 \(f\) 在 \(\mathrm{L}^2 (\mathrm{G})\) 中是族 \((f_{\overline{\pi}})\) 的和，根据彼得–外尔定理，得到

\[
f = \sum_ {\pi \in \widehat {\mathrm{G}}} \mathcal {F} _ {\pi} (\pi (f))
\]

其中级数在 \(L^{2}(G)\) 中收敛。这就证明了定理。

定义 2. --- \(\mathcal{F}_{\mathrm{G}}\) 是从 \(\mathrm{F}^2 (\widehat{\mathrm{G}})\) 到 \(\mathrm{L}^2 (\mathrm{G})\) 的等距同构，也是傅里叶余变换所诱导同构的逆；称其为 G 的傅里叶变换。称 \(\mathrm{F}^2 (\widehat{\mathrm{G}})\) 中元素的像为该元素的傅里叶变换。

注记。--- 1) \((u_{\pi})_{\pi \in \widehat{\mathrm{G}}} \in \mathrm{F}^2(\widehat{\mathrm{G}})\) 的傅里叶变换因此是下列级数在 \(\mathrm{L}^2(\mathrm{G})\) 中的等价类

\[
g \mapsto \sum_ {\pi \in \widehat {\mathrm{G}}} \langle \pi (g) | u _ {\pi} \rangle = \sum_ {\pi \in \widehat {\mathrm{G}}} \dim (\pi) \operatorname{Tr} (u _ {\pi} \circ \pi (g) ^ {- 1}).
\]

\begin{enumerate}
\def\labelenumi{\arabic{enumi})}
\setcounter{enumi}{1}
\tightlist
\item
  令 \(f \in \mathrm{L}^2(\mathrm{G})\)。于是有普朗歇雷尔公式
\end{enumerate}

\[
\| f \| ^ {2} = \| \overline {{\mathcal {F}}} _ {\mathrm{G}} (f) \| ^ {2} = \sum_ {\pi \in \widehat {\mathrm{G}}} \| \pi (f) \| ^ {2}.
\]

此外，根据定理 2，\(f\) 是 \(\mathrm{L}^2 (\mathrm{G})\) 中族 \((f_{\pi})_{\pi \in \widehat{\mathrm{G}}}\) 的和，其中

\[
\begin{array}{l} f _ {\pi} (x) = \mathcal {F} _ {\pi} (\pi (f)) (x) \\ \qquad = \langle \pi (x) | \pi (f) \rangle = \dim (\pi) \int_ {\mathrm{G}} f (g) \chi_ {\pi} (x ^ {- 1} g) d \mu (g) \end{array}
\]

对所有 \(x \in G\) 及每个 \(\pi \in \widehat{G}\) 都成立。

设 G 为交换群。由于 G 的不可约表示维数为 1（V, 第 390 页推论 7），且对偶群 \(\widehat{\mathbf{G}}\) 是离散的（II, 第 233 页命题 18），代数 \(\mathrm{F}_b(\widehat{\mathbf{G}})\) 和 \(\mathrm{F}_0(\widehat{\mathbf{G}})\) 分别识别为代数 \(\mathcal{C}_b(\widehat{\mathbf{G}})\)（在 \(\widehat{\mathbf{G}}\) 上的连续有界函数的代数）和代数 \(\mathcal{C}_0(\widehat{\mathbf{G}})\)（在 \(\widehat{\mathbf{G}}\) 上在无穷远处趋于 0 的连续函数的代数）。由于 G 是紧的，\(\widehat{\mathbf{G}}\) 上与 \(\mu\) 对偶的哈尔测度是计数测度 \(\widehat{\mu}\)（II, 第 233 页命题 18）。希尔伯特直和 \(\mathrm{F}^2 (\widehat{\mathrm{G}})\) 由各空间 \(\operatorname{End}(\mathrm{E}_{\pi})\)（\(\pi \in \widehat{\mathrm{G}}\)）组成，并识别为希尔伯特空间 \(\mathrm{L}^2 (\widehat{\mathrm{G}},\widehat{\mu})\)。

令 \(\nu \in \mathcal{M}^1 (\mathrm{G})\)。于是对每个酉特征标 \(\chi \in \widehat{\mathrm{G}}\)，有

\[
\chi (\nu) = \int_ {\mathrm{G}} \chi \nu
\]

这证明了上面定义的傅里叶余变换与 II, 第 206 页定义的 G 的傅里叶余变换一致。

每个 \(f \in \mathcal{L}^2(\mathrm{G})\) 都可积，而公式 \(\| \overline{\mathcal{F}}_{\mathrm{G}}(f)\| _2 = \| f\|\) 就是普朗歇雷尔公式（II, 第 215 页定理 1）。

命题 10. --- 傅里叶余变换是从 \(\mathcal{M}^{1}(\mathrm{G})\) 到 \(\mathrm{F}_{b}(\widehat{\mathrm{G}})\) 的单射对合巴拿赫代数态射，并将 \(\mathrm{L}^{1}(\mathrm{G})\) 映入 \(\mathrm{F}_{0}(\widehat{\mathrm{G}})\)。

由于 G 是紧的，空间 \(\mathcal{C}(\mathrm{G})\) 包含于 \(\mathcal{L}^{1}(\mathrm{G})\) 和 \(\mathcal{L}^{2}(\mathrm{G})\)，并且在二者中都稠密。

令 \(\nu \in \mathcal{M}^1 (\mathrm{G})\) 满足 \(\overline{\mathcal{F}}_{\mathrm{G}}(\nu) = 0\)。于是对 G 上的每个连续函数 \(f\)，都有 \(\overline{\mathcal{F}}_{\mathrm{G}}(\nu *f) = 0\)。现在 \(\nu *f\) 属于 \(\mathcal{C}(\mathrm{G})\)（INT, VIII, 第 152 页，第 4 节第 2 号命题 3）。由于根据定理 2，傅里叶余变换在 \(\mathrm{L}^2 (\mathrm{G})\) 上是单射，因而在空间 \(\mathcal{C}(\mathrm{G})\) 上也是单射，所以 \(\nu *f = 0\) 对 \(f\in \mathcal{C}(\mathrm{G})\) 成立。特别地，可得

\[
\int_ {\mathrm{G}} f (x) d \nu (x) = (\nu * \check {f}) (e) = 0
\]

对每个函数 \(f \in \mathcal{C}(\mathrm{G})\) 都成立（INT, VIII, 上引处），所以测度 \(\nu\) 为零。

傅里叶余变换下 \(\mathcal{C}(\mathrm{G})\) 的像包含于 \(\mathrm{F}^2 (\widehat{\mathrm{G}})\)，从而也包含于 \(\mathrm{F}_0(\widehat{\mathrm{G}})\)。由于 \(\mathrm{F}_0(\widehat{\mathrm{G}})\) 在 \(\mathrm{F}_b(\widehat{\mathrm{G}})\) 中是闭的，且 \(\mathcal{C}(\mathrm{G})\) 在 \(\mathrm{L}^1 (\mathrm{G})\) 中稠密，傅里叶余变换下 \(\mathrm{L}^1 (\mathrm{G})\) 的像也包含于 \(\mathrm{F}_0(\widehat{\mathrm{G}})\)。

命题 11. --- a) 令 \(u = (u_{\pi})_{\pi \in \widehat{\mathbf{G}}}\) 为 \(\mathrm{F}(\widehat{\mathbf{G}})\) 中的元素。若族 \((\mathcal{F}_{\pi}(u_{\pi}))_{\pi \in \widehat{\mathbf{G}}}\) 在 \(\mathcal{C}(\mathbf{G})\) 中一致可求和，则其和 \(f\) 是 \(\mathbf{G}\) 上的连续函数，且其傅里叶余变换为 \(u\)；

\begin{enumerate}
\def\labelenumi{\alph{enumi})}
\setcounter{enumi}{1}
\tightlist
\item
  令 \(f \in \mathrm{L}^1(\mathrm{G})\)。若族 \((\mathcal{F}_\pi(\pi(f)))_{\pi \in \widehat{\mathrm{G}}}\) 在 \(\mathcal{C}(\mathrm{G})\) 中一致可求和，则其和是 \(\mathrm{G}\) 上的连续函数，且它在 \(\mathrm{L}^1(\mathrm{G})\) 中的等价类等于 \(f\)。
\end{enumerate}

证明断言 \(a\)。由于 G 是紧的，\(f\) 是族 \((\mathcal{F}_{\pi}(u_{\pi}))\) 的和，属于 \(\mathrm{L}^2 (\mathrm{G})\)，且级数

\[
\sum_ {\pi \in \widehat {\mathrm{G}}} \mathcal {F} _ {\pi} (u _ {\pi})
\]

在 \(L^{2}(G)\) 中收敛于 f（INT, IV, 第 127 页，第 3 节第 3 号命题 4）。因此 \(\mathcal{F}_{\mathrm{G}}((u_{\pi})_{\pi}) = f\) 在 \(L^{2}(G)\) 中成立，从而 \((u_{\pi})_{\pi \in \widehat{\mathrm{G}}} = \overline{\mathcal{F}}_{\mathrm{G}}(f)\)（定理 2）。

证明断言 b)。可以将断言 a) 应用于族 \((\pi(f))_{\pi \in \widehat{\mathrm{G}}}\)。族 \((\mathcal{F}_{\pi}(\pi(f)))_{\pi \in \widehat{\mathrm{G}}}\) 的和 g 是连续函数，因而属于 \(L^{2}(G)\)，并满足 \(\overline{\mathcal{F}}_{\mathrm{G}}(g) = (\pi(f))_{\pi \in \widehat{\mathrm{G}}}\)。由于 \(\overline{\mathcal{F}}_{\mathrm{G}}\) 在 \(L^{2}(G)\) 上是单射，所以在 \(L^{2}(G)\) 中 f = g，因而在 \(L^{1}(G)\) 中也有 f = g。

将 \(\mathcal{C}(\widehat{\mathrm{G}})\) 识别为 \(\widehat{\mathrm{G}}\) 上复值函数的代数，它是代数 \(\mathrm{F}(\widehat{\mathrm{G}})\) 的中心：该映射将 \(f\colon \widehat{\mathrm{G}}\to \mathbf{C}\) 对应于中心元素 \((f(\pi)1_{\mathrm{E}_{\pi}})_{\pi \in \widehat{\mathrm{G}}}\)（属于 \(\mathrm{F}(\widehat{\mathrm{G}})\)）。记 \(\mathcal{M}^1 (\mathrm{G}^\sharp)\) 为 G 上中心有界测度的空间。它是巴拿赫代数 \(\mathcal{M}^1 (\mathrm{G})\) 的闭子空间。对于每个测度 \(\nu \in \mathcal{M}^1 (\mathrm{G}^\sharp)\) 和每个表示 \(\pi \in \widehat{\mathrm{G}}\)，有 \(\pi (\nu)\in \mathrm{End}_{\mathrm{G}}(\mathrm{E}_{\pi})\)，因而根据舒尔引理（V, 第 386 页命题 6），\(\pi (\nu)\) 是 \(\mathrm{E}_{\pi}\) 的恒等自同态的数量倍。傅里叶余变换在 \(\mathcal{M}^1 (\mathrm{G}^\sharp)\) 上的限制因而识别为从 \(\mathcal{M}^1 (\mathrm{G}^\sharp)\) 到 \(\mathcal{C}(\widehat{\mathrm{G}})\) 的含单位对合巴拿赫代数态射。该态射是单射；\(\mathrm{L}^1 (\mathrm{G}^\sharp)\) 的像包含于 \(\mathcal{C}_b(\widehat{\mathrm{G}})\)，而其在 \(\mathrm{L}^2 (\mathrm{G}^\sharp)\) 上的限制是到 \(\mathrm{L}^2 (\widehat{\mathrm{G}},\widehat{\mu})\) 的等距同构。

\subsection{弗罗贝尼乌斯--舒尔指标与拉森择一定理}

回忆（LIE, VIII, 第 7 节第 6 号定义 2）：若 G 在有限维复向量空间 E 中的不可约表示存在一个在 G 下不变的非零对称双线性形式（分别地，存在一个在 G 下不变的非零交错双线性形式），则称该表示为正交型（分别为辛型）。若 E 上不存在在 G 下不变的非零双线性形式，则称其为复型。

当 E 中的不可约表示为正交型（分别为辛型）时，E 上 G-不变对称（分别为交错）双线性形式的空间是一维的，并且每个非零 G-不变双线性形式都是分离的。

正交型表示有时称为实型表示，辛型表示有时称为四元数型表示（参见 LIE, IX, 附录 II）。

令 \(\pi\) 为 G 在复向量空间 E 中的不可约表示。上述三种情形由量

\[
\mathrm{FS} (\pi) = \int_ {\mathrm{G}} \chi_ {\pi} (g ^ {2}) d \mu (g),
\]

区分；称其为 π 的弗罗贝尼乌斯--舒尔指标。事实上有

\[
\operatorname{FS} (\pi) = \left\{ \begin{array}{l l} 1 & \text{若 }\pi\text{ 属于正交型，} \\ - 1 & \text{若 }\pi\text{ 属于辛型，} \\ 0 & \text{若 }\pi\text{ 属于复型。} \end{array} \right.
\]

（LIE, IX, 附录 2，第 103 页命题 3 及第 105 页命题 4）。

当 G 为李群时，可以利用 LIE, IX, 第 69 页命题 1 计算 \(\mathrm{FS}(\pi)\)。

定义 3. --- 令 \(\varrho\) 为 G 在希尔伯特空间 E 中的有限维酉表示。令 \(k\) 为正整数。称表示 \(\overline{\varrho}^{\otimes k} \otimes \varrho^{\otimes k}\) 的 G-不变元素子空间的维数为 \(2k\) 阶的 \(\varrho\) 的绝对矩，记为 \(\mathrm{M}_{2k}(\varrho)\)。

定理 3（拉森择一定理）。--- 设 G 是无限群。令 π 为 G 在维数 ≥ 2 的有限维希尔伯特空间 E 中的忠实酉表示。

\begin{enumerate}
\def\labelenumi{\alph{enumi})}
\item
  假设 G 的导出群是无限的。有 \(M_{4}(\pi) \geqslant 2\)，且等号成立当且仅当 G 包含 SU(E)；
\item
  假设 \(\dim(E) \geqslant 3\)，\(\pi\) 为正交型或辛型，且 \(\dim(E) \neq 4\)（在 \(\pi\) 为正交型时）。则 \(M_4(\pi) \geqslant 3\)，且等号成立当且仅当 G 的复化李代数等于 E 上某个分离的 G-不变双线性形式的李代数。*当且仅当 G 的单位分支 \(G_0\) 是这种双线性形式的自同构群的极大紧子群时，情形才成立。*
\end{enumerate}

证明中要用到下列引理。

引理 5. --- 令 \(\pi\) 为 G 在有限维希尔伯特空间 E 中的酉表示。令 \((\varrho_{i})_{i\in\mathrm{I}}\) 为 G 的非零酉表示族，\((n_{i})_{i\in\mathrm{I}}\) 为满足 \(\geqslant 1\) 的整数族，使得 G 的表示 \(\overline{\pi}\otimes\pi\)（分别地，表示 \(\pi\otimes\pi\)）同构于直和 \(\bigoplus_{i\in\mathrm{I}}\varrho_{i}^{n_{i}}\)。则有

\[
\mathrm{M} _ {4} (\pi) \geqslant \sum_ {i \in \mathrm{I}} n _ {i} ^ {2}
\]

等号成立当且仅当表示 \(\varrho_{i}\) 都不可约且两两不同构。

记 \(\chi_{i}\) 为 \(\varrho_{i}\) 的特征标，其中 \(i\in I\)。有

\[
| \chi_ {\pi} | ^ {2} = \chi_ {\overline {{\pi}} \otimes \pi} = \sum_ {i \in \mathrm{I}} n _ {i} \chi_ {i}, \quad (\text { resp. } \chi_ {\pi} ^ {2} = \chi_ {\pi \otimes \pi} = \sum_ {i \in \mathrm{I}} n _ {i} \chi_ {i}).
\]

根据定义和 V, 第 466 页推论 3，有

\[
\mathrm{M} _ {4} (\pi) = \int_ {\mathrm{G}} | \chi_ {\pi} | ^ {4} d \mu ,
\]

从而在两种情形下均有公式

\[
\mathrm{M} _ {4} (\pi) = \sum_ {i, j} n _ {i} n _ {j} \int_ {\mathrm{G}} \chi_ {i} \overline {{\chi}} _ {j} d \mu = \sum_ {i, j} n _ {i} n _ {j} \dim \mathrm{Hom} _ {\mathrm{G}} (\varrho_ {i}, \varrho_ {j})
\]

（V, 第 459 页推论 4, b)）。令 \(i \in I\)。由于表示 \(\varrho_{i}\) 非零，空间 \(\mathrm{Hom}_{\mathrm{G}}(\varrho_{i}, \varrho_{i})\) 非零，于是得到下界

\[
\mathrm{M} _ {4} (\pi) \geqslant \sum_ {i \in \mathrm{I}} n _ {i} ^ {2}.
\]

此外，由于 \(n_{i} \geqslant 1\)，等号成立当且仅当 \(\dim \operatorname{Hom}_{\mathrm{G}}(\varrho_{i}, \varrho_{j}) = 0\)（当 \(i \neq j\) 时），并且对每个 i 都有 \(\dim \operatorname{Hom}_{\mathrm{G}}(\varrho_{i}, \varrho_{i}) = 1\)。第二个条件成立当且仅当表示 \(\varrho_{i}\) 不可约（上引处）。此时根据舒尔引理（V, 第 387 页推论 2），第一个条件成立当且仅当表示 \(\varrho_{i}\) 两两不同构。

若 E 是交换环 A 上的模，则称 E 的对称平方（分别为外平方）为 A-模 \(S^{2}(E)\)（分别为 \(\wedge^{2}E\)）。

引理 6. --- 令 \(q\) 为维数 \(\geqslant 3\) 的有限维复向量空间 E 上的分离双线性形式。

\begin{enumerate}
\def\labelenumi{\alph{enumi})}
\item
  若 \(q\) 是对称的，则正交李代数 \(\mathfrak{so}(q)\) 的伴随表示同构于 \(\wedge^2\mathrm{E}\)，即自然表示 \(\mathfrak{so}(q) \to \mathfrak{gl}(\mathrm{E})\) 的外平方；
\item
  若 q 是交错的，则辛李代数 \(\mathfrak{sp}(q)\) 的伴随表示同构于 \(\mathsf{S}^2 (\mathrm{E})\)，即自然表示 \(\mathfrak{sp}(q)\to \mathfrak{gl}(\mathrm{E})\) 的对称平方。
\end{enumerate}

在 C 上赋予恒等对合自同构。于是双线性形式 q 是 \(\varepsilon\)-埃尔米特的（A, IX, 第 3 节第 1 号定义 1），其中情形 a) 中 \(\varepsilon = 1\)，情形 b) 中 \(\varepsilon = -1\)。对于每个 \(u \in \text{End}(\text{E})\)，记 \({}^{t}u\) 为关于 q 的 u 的伴随。于是有 \(q(x, u(y)) = q(^{t}u(x), y)\)，对所有 \((x, y) \in \text{E} \times \text{E}\) 成立。令 v 为 E \(\otimes\) E 的唯一自同构，使得 \(v(x \otimes y) = y \otimes x\) 对所有 \((x, y) \in \text{E}^{2}\) 成立。

存在唯一的线性映射 \(w\) 从 \(\mathrm{E} \otimes \mathrm{E}\) 到 \(\operatorname{End}(\mathrm{E})\)，使得 \(w(a \otimes b)\) 是由 \(x \mapsto q(a, x)b\) 定义的线性映射，对所有 \((a, b, x) \in \mathrm{E}^3\) 成立。映射 \(w\) 是同构。对于每个 \(s \in \mathrm{E} \otimes \mathrm{E}\)，有

\[
{ } ^ { t } ( w ( s ) ) = \varepsilon w ( v ( s ) ) .\tag{3}
\]

事实上，对于 \((x,y)\) 和 \((a,b)\) 在 \(\mathrm{E} \times \mathrm{E}\) 中的所有情形，可得

\[
q (x, w (a \otimes b) y) = q (a, y) q (x, b) = \varepsilon q (b, x) q (a, y) = q (\varepsilon w (b \otimes a) x, y).
\]

最后，记 \(x \mapsto x^*\) 为由 q 导出的从 E 到 \(\mathrm{E}'\) 的同构，即由 \(q\) 给出，满足 \(\langle x^*, y \rangle = q(x, y)\)，对所有 \((x, y) \in \mathrm{E}^2\) 成立。

记号确定后，令 H 为 q 的正交群（分别为辛群）。它是 \(\mathbf{GL}(\mathrm{E})\) 的李子群，其李代数 h 是由 \(u \in \operatorname{End}(\operatorname{E})\) 满足 \(^{t}u = -u\) 的 End(E) 子空间（LIE, III, 第 146 页推论 1）。由于形式 q 在 H 下不变，故有

\[
\langle (g a) ^ {*}, b \rangle = q (g a, b) = \langle a, g ^ {- 1} b \rangle
\]

对所有 \(g \in H\) 及每个 \((a, b) \in E \times E\) 都成立。

在 End(E) 上赋予 H 的伴随表示 Ad。线性映射 \(w\) 是 H 的表示之间的态射：事实上，对每个 \(g \in \mathrm{H}\) 及每个 \((a, b, x) \in \mathrm{E}^3\)，有

\[
\begin{array}{c} w (g   (a \otimes b)) (x) = \langle (g a) ^ {*}, x \rangle   g b = \langle a, g ^ {- 1} x \rangle   g b \\ = g   (\langle a, g ^ {- 1} x \rangle   b) = \mathrm{Ad} (g) (w (a \otimes b)) x, \end{array}
\]

由线性性即得结论。

由于 v 也是 H 的表示之间的态射，线性映射 \(\theta = w - \varepsilon w \circ v\) 也具有这一性质；因此它是 h-模之间的态射。

记 \(F \subset E \otimes E\) 为由元素 \(s \in E \otimes E\) 且满足 \(v(s) = -\varepsilon s\) 所成的子空间。若 \(\varepsilon = -1\)，则从 \(E \otimes E\) 到 E 的对称平方的典范映射在 F 上的限制是 C-向量空间同构（A, III, 第 72 页）；若 \(\varepsilon = 1\)，则从 \(E \otimes E\) 到外平方的典范映射在 F 上的限制是 C-向量空间同构（上引处，第 82 页）。

F 在 \(\theta\) 下的像包含于 \(\mathfrak{h}\) 中：事实上，对每个 \(s\in\mathrm{F}\)，公式 (3) 蕴含

\[
{ } ^ { t } \theta ( s ) = { } ^ { t } w ( s ) - \varepsilon { } ^ { t } w ( v ( s ) ) = \varepsilon w ( v ( s ) ) - w ( s ) = - \theta ( s ) .
\]

按 F 的定义，映射 \(\theta\) 在 F 上的限制与 \(2w\) 的限制重合，因而线性映射 \(\theta\) 在 F 上是单射。由于 \(\dim(F) = \dim(h)\)。由（A, III, 第 75 页定理 1 和第 87 页推论 1，以及 LIE, VIII, 第 192 页和第 201 页），可知 \(\theta\) 通过取子空间诱导出从 F 到 \(h\) 的 \(h\)-模同构，故引理得证。

引理 7. --- 令 \(H_{2}\) 为实李群，令 \(H_{1}\) 为 \(H_{2}\) 的闭子群。\(H_{1}\) 的复化李代数可识别为 \(H_{1}\) 在 \(H_{2}\) 的复化李代数上的伴随表示的一个子表示。

事实上，限制到 \(H_{1}\) 的 \(H_{2}\) 李代数上的伴随表示，正可识别为 \(H_{1}\) 的伴随表示。

引理 8. --- 令 E 为维数 n 的复希尔伯特空间。

\begin{enumerate}
\def\labelenumi{\alph{enumi})}
\item
  群 \(\mathbf{SU}(\mathrm{E})\) 和 \(\mathbf{U}(\mathrm{E})\) 是连通的；
\item
  SU(E) 的导出群等于 SU(E)；
\item
  U(E) 的导出子群等于 SU(E)。
\end{enumerate}

不妨设 \(n \geqslant 1\) 且 \(E = C^{n}\)。

令 A 为 \(\mathbf{U}(\mathrm{E})\) 中由对角矩阵组成的子群；它同胚于 \(U^{n}\)，因而是连通的（TG, VI, 第 11 页推论 2 以及 I, 第 83 页命题 8）。它与 \(\mathbf{SU}(\mathrm{E})\) 的交同胚于 \(U^{n-1}\)，所以也是连通的。

由 IV, 第 149 页定理 1，群 \(\mathbf{U}(\mathrm{E})\)（分别为 \(\mathbf{SU}(\mathrm{E})\)）是连通子集 \(g\mathrm{Ag}^{-1}\)（\(g \in \mathrm{G}\)）的并（分别是连通子集 \(g(\mathrm{A} \cap \mathbf{SU}(\mathrm{E}))g^{-1}\) 的并）；由于单位元属于这些集合中的每一个，空间 \(\mathbf{U}(\mathrm{E})\)（分别为 \(\mathbf{SU}(\mathrm{E})\)）是连通的（TG, I, 第 81 页命题 2）。

证明 \(b\)。当 \(n = 1\) 时断言成立，因为此时 \(\mathbf{SU}(\mathrm{E})\) 退化为单位元。因此设 \(n \geqslant 2\)。此时 \(\mathbf{SU}(\mathrm{E})\) 的李代数是单的（LIE, IX, 第 20 页，第 3 节第 4 号），所以 \(\mathbf{SU}(\mathrm{E})\) 的导出群在 \(\mathbf{SU}(\mathrm{E})\) 中具有有限指数。由于 \(\mathbf{SU}(\mathrm{E})\) 由 \(a\) 可知连通，故得到所需结果。

断言 \(c\) 由 \(b\) 得出，因为 \(\mathbf{U}(\mathrm{E})\) 的导出群包含于 \(\mathbf{SU}(\mathrm{E})\)。

现在证明定理 3。记 n 为希尔伯特空间 E 的维数。由于表示 \(\pi\) 忠实，不妨设 G 是 U(E) 的紧子群。群 G 在实李群 U(E) 中是闭的，因而是紧实李群（LIE, III, 第 8 节第 2 号定理 2）。根据假设，G 是无限的，故其维数 \(\geqslant 1\)。令 g 为其李代数；它非零。

证明 \(a\)。假设导出群 D(G) 是无限的。有 G-不变分解 \(\operatorname{End}(\mathrm{E}) = \mathbf{C}1_{\mathrm{E}} \oplus \mathfrak{sl}(\mathrm{E})\)。G 在 \(\mathbf{C}1_{\mathrm{E}}\) 和 \(\mathfrak{sl}(\mathrm{E})\) 上的表示不同构，因为 \(\dim \mathfrak{sl}(\mathrm{E}) \geqslant 2\)。由引理 5，\(\mathrm{M}_4(\pi) \geqslant 2\)，且等号成立当且仅当 G 在 \(\mathfrak{sl}(\mathrm{E})\) 中的表示不可约。现在，G 的导出群包含于 \(\mathbf{SL}(\mathrm{E})\)，所以 D(G) 的李代数的复化可识别为 \(\mathfrak{sl}(\mathrm{E})\) 的一个子空间，而后者是 G 在 \(\mathfrak{sl}(\mathrm{E})\) 上的表示的一个子表示（引理 7）。由于 D(G) 无限，该子表示非零。因此，\(\mathrm{M}_4(\pi) = 2\) 当且仅当 \((\mathcal{D}\mathfrak{g})_{(\mathbf{C})} = \mathfrak{sl}(\mathrm{E})\)。由于 D(G) 包含于 \(\mathbf{SU}(\mathrm{E})\)，且 \(\mathbf{SU}(\mathrm{E}) = \mathrm{D}(\mathbf{SU}(\mathrm{E}))\)（引理 8），这个条件等价于 \(\mathbf{SU}(\mathrm{E}) \subset \mathbf{G}\)。

为证明断言 \(b\)，沿用 LIE, VIII, 第 13 节第 2--4 号的记号。

设 \(\pi\) 为辛型且 \(\dim(\mathrm{E}) \geqslant 3\)。令 \(q\) 为 E 上非零的 G-不变交错双线性形式；它是分离的（LIE, IX, 第 103 页附录 2 命题 3），且 G 是辛群 \(\mathbf{Sp}(q)\) 的紧子群。记 \(q^{*} \in \wedge^{2}\mathrm{E}\) 为交错形式 \(q\) 所对应的元素。G 的复化李代数 \(\mathfrak{g}_{(\mathbf{C})}\) 包含于 \(\mathfrak{sp}(q)\)。现在，由于 \(\dim(\mathrm{E}) \geqslant 3\)，\(\mathfrak{sp}(q)\) 在 \(\mathrm{E} \otimes \mathrm{E}\) 中的表示具有如下直和分解

\[
\mathrm{E} \otimes \mathrm{E} = \mathsf {S} ^ {2} (\mathrm{E}) \oplus \wedge^ {2} \mathrm{E} = \mathsf {S} ^ {2} (\mathrm{E}) \oplus \mathrm{E} _ {2} \oplus \mathbf {C} q ^ {*},
\]

其中 \(E_{2}\) 是 \(\mathfrak{sp}(q)\) 的第二基本表示（LIE, VIII, 第 202 页，第 13 节第 3 号 (IV)）。表示 \(E_{2}\) 和 \(C q^{*}\) 是 \(\mathfrak{sp}(q)\) 的不可约表示（上引处），而表示 \(S^{2}(E)\) 非零。

因此由引理 5，\(\mathrm{M}_{4}(\pi) \geqslant 3\)，且等号成立当且仅当 G 在 \(\mathsf{S}^{2}(\mathrm{E})\)、\(E_{2}\) 和 \(Cq^{*}\) 中的表示都不可约且两两不同构。

设 \(M_{4}(\pi)=3\)。由于表示 \(S^{2}(E)\) 可识别为 \(\mathfrak{sp}(q)\) 的伴随表示（引理 6），它包含 G 的伴随表示的复化作为子表示（引理 7）。因此有 \(\mathfrak{g}_{(\mathbf{C})}=\mathfrak{sp}(q)\)。

反之，设 \(\mathfrak{g}_{(\mathbf{C})} = \mathfrak{sp}(q)\)。\(\mathfrak{sp}(q)\) 的伴随表示和表示 \(\mathrm{E}_2\)，后者属于 \(\mathfrak{sp}(q)\)，都不可约，其维数分别为 \(n(n + 1)/2\) 和 \(n(n - 1)/2 - 1\)（LIE, VIII, 第 202 页，第 13 节第 3 号 (IV)）；这两个维数不同且均 \(\geqslant 2\)，因为 \(n \geqslant 3\)，从而 \(\mathrm{M}_4(\pi) = 3\)。

最后，设 \(\pi\) 为实型，且 \(\dim(\mathrm{E}) \geqslant 3\) 以及 \(\dim(\mathrm{E}) \neq 4\)。令 \(q\) 为 \(\mathrm{E}\) 上的非零对称双线性形式，且在 G 下不变；它是分离的（LIE, IX, 第 103 页附录 2 命题 3），并且 \(\mathrm{G}\) 是正交群 \(\mathbf{O}(q)\) 的紧子群。记 \(q^{*} \in \mathbb{S}^{2}(\mathrm{E})\) 为与 \(q\) 对应的元素。

G 的复化李代数 \(\mathfrak{g}_{(\mathbf{C})}\) 包含于 \(\mathfrak{so}(q)\)。\(\mathfrak{so}(q)\) 在 E \(\otimes\) E 中的表示具有如下直和分解

\[
\mathrm{E} \otimes \mathrm{E} = \wedge^ {2} \mathrm{E} \oplus \mathsf {S} ^ {2} (\mathrm{E}) = \wedge^ {2} \mathrm{E} \oplus \mathsf {S} _ {0} ^ {2} (\mathrm{E}) \oplus \mathbf {C} q ^ {*},
\]

其中 \(\mathsf{S}_{0}^{2}(\mathbf{E})\) 是 \(q^{*}\) 在 \(\mathsf{S}^{2}(\mathbf{E})\) 中的正交补。这些表示的维数至少为 2，因为 \(\dim(\mathbf{E})\geqslant3\)。由引理 5，\(\mathrm{M}_{4}(\pi)\geqslant3\)，且等号成立当且仅当 G 在 \(\wedge^{2}\mathrm{E}\) 和 \(\mathsf{S}_{0}^{2}(\mathbf{E})\) 中的表示不可约且不同构。

设 \(\mathrm{M}_{4}(\pi)=3\)。由于表示 \(\wedge^{2}E\) 可识别为 \(\mathfrak{so}(q)\) 的伴随表示（引理 6），它包含 G 的伴随表示的复化作为子表示（引理 7）。因此条件 \(\mathrm{M}_{4}(\pi)=3\) 蕴含 \(\mathfrak{g}_{(\mathbf{C})}=\mathfrak{so}(q)\)。

反之，设 \(\mathfrak{g}_{(\mathbf{C})} = \mathfrak{so}(q)\)。\(\mathfrak{so}(q)\) 的伴随表示是不可约的，因为 \(\dim(\mathrm{E}) \neq 4\)（LIE, VIII, 第 13 节，第 193 页 (I) 及第 206 页 (I)），且维数为 \(n(n-1)/2\)。表示 \(\mathsf{S}_0^2(\mathrm{E})\) 是 \(\mathfrak{so}(q)\) 的表示，其最高权为 \(2\varpi_1\)。将其维数与由 Weyl 公式计算的 \(\mathfrak{so}(q)\) 的最高权为 \(2\varpi_1\) 的不可约表示的维数比较（参见 LIE, VIII, 第 9 节第 2 号定理 2 以及 LIE, VI, 图版 II 和 IV），可验证 \(\mathsf{S}_0^2(\mathrm{E})\) 是不可约的。因此得到 \(\mathrm{M}_4(\pi) = 3\)，当 \(\mathfrak{g}_{(\mathbf{C})} = \mathfrak{so}(q)\) 时。

注记。--- 1) “G 是无限的”这一条件对于定理 3 是必要的（V, 第 507 页习题 9）。

\begin{enumerate}
\def\labelenumi{\arabic{enumi})}
\setcounter{enumi}{1}
\item
  令 \(n \in \mathbf{N}\)。对于紧群的每个酉表示 \(\varrho\)，若它作用在维数为 \(n\) 的希尔伯特空间中，则有 \(\mathrm{M}_{4}(\varrho) \leqslant n^{2}\)；等号可能成立（V, 第 508 页习题 15）。
\item
  可以证明（参见 R. GURALNICK 和 P.H. TIEP，《Decomposition of small tensor powers and Larsen's conjecture》，Representation Theory 9 (2005), 138--208）：若假设 E 的维数 \(\geqslant 7\)，并将假设 \(M_{4}(\pi) = 2\)（分别为 \(M_{4}(\pi) = 3\)）替换为 \(M_{8}(\pi) = 24\)（分别为 \(M_{8}(\pi) = 105\)），则可以省略 G 为无限这一条件。
\end{enumerate}

\specialchapter{习题}

\exercisesection{}

\begin{enumerate}
\def\labelenumi{\arabic{enumi})}
\item
  给出一个群 G、一个域 K 以及 G 的一个忠实表示的例子，使得相应的 \(K[G]\)-模并不忠实。
\item
  令 \(\pi\) 为 G 在希尔伯特空间 E 上的酉表示，且 \(x \in E\) 的范数为 1，并满足
\end{enumerate}

\[
\sup _ {g \in \mathrm{G}} \| \pi (g) x - x \| \leqslant 1.
\]

存在非零的 G-不变向量 \(y \in E\)。

\begin{enumerate}
\def\labelenumi{\arabic{enumi})}
\setcounter{enumi}{2}
\tightlist
\item
  令 E 为希尔伯特空间。记 \(\mathbf{U}(\mathbf{E})\) 为 E 的酉自同态群，记 \(\mathbf{GL}(\mathbf{E})\) 为 E 的可逆自同态群。
\end{enumerate}

\begin{enumerate}
\def\labelenumi{\alph{enumi})}
\item
  群 \(\mathbf{U}(\mathrm{E})\) 赋予由 \(\mathcal{L}(\mathrm{E})\) 上逐点收敛拓扑诱导的拓扑后是拓扑群，记为 \(\mathbf{U}(\mathrm{E})_{s}\)；若 E 是无限维的，则 \(\mathbf{GL}(\mathrm{E})\) 赋予由逐点收敛拓扑诱导的拓扑后不是拓扑群。
\item
  \(\mathbf{U}(\mathrm{E})_{s}\) 的拓扑与由紧收敛拓扑诱导的拓扑一致。
\item
  若 E 是可数生成的，则群 \(\mathbf{U}(\mathrm{E})_{s}\) 可度量化且完备。
\item
  令 L 为由 \(u \in \mathcal{L}(\mathrm{E})\) 且满足 \(u^{*} = -u\) 的元素所成的空间；它是实李群 G = U(E) 的李代数，其中 G 赋予由 \(\mathcal{L}(\mathrm{E})\) 的巴拿赫空间拓扑诱导的拓扑（LIE, III, 第 146 页推论 2）。从 L 到 \(\mathbf{U}(\mathrm{E})_s\) 的指数映射是连续的。
\item
  若 E 是无限维的，则 \(\mathbf{U}(\mathrm{E})_{s}\) 不是型 \(\{L\}\) 的实李群。（考虑 \(\mathbf{U}(\mathrm{E})_{s}\) 中由 E 的某个固定标准正交基下的对角酉自同态组成的子群 K，并注意 K 是紧的。）
\item
  指数映射不是从 L 到 \(\mathbf{U}(\mathbf{E})_{s}\) 的局部同胚。（考虑它在某个固定标准正交基下的对角自同态上的限制。）
\end{enumerate}

\begin{enumerate}
\def\labelenumi{\arabic{enumi})}
\setcounter{enumi}{3}
\tightlist
\item
  令 G 为拓扑群。若存在一个线性泛函 \(\lambda\colon\mathcal{C}_{b}(\mathrm{G};\mathbf{R})\to\mathbf{R}\)，使得
\end{enumerate}

\begin{enumerate}
\def\labelenumi{(\roman{enumi})}
\item
  若 \(f \geqslant 0\)，则 \(\lambda(f) \geqslant 0\)。
\item
  有 \(\lambda(1)=1\)。
\item
  对每个 \(g \in \mathbf{G}\) 及每个 \(f \in \mathcal{C}_b(\mathbf{G};\mathbf{R})\)，有 \(\lambda(\boldsymbol{\gamma}(g)f) = \lambda(\boldsymbol{\delta}(g)f)\)，其中
\end{enumerate}

\[
\boldsymbol {\gamma} (g) f (x) = f \left(g ^ {- 1} x\right), \quad \boldsymbol {\delta} (g) f (x) = f (x g).
\]

若 G 是离散群，此定义与 EVT, IV, 第 73 页习题 4 中的定义一致。

\begin{enumerate}
\def\labelenumi{\alph{enumi})}
\tightlist
\item
  拓扑群 G 可和，当且仅当对于每个整数 \(n \in N\)、每个有限族 \((f_i)_{1 \leqslant i \leqslant n}\)（属于 \(\mathcal{C}_b(\mathbf{G}; \mathbf{R})\)）、G 中的两个族 \((g_i)_{1 \leqslant i \leqslant n}\) 和 \((h_i)_{1 \leqslant i \leqslant n}\) 以及每个实数 a，条件
\end{enumerate}

\[
\sum_ {i = 1} ^ {n} (f _ {i} (x) - f _ {i} (g _ {i} x h _ {i})) \geqslant a
\]

对所有 \(x \in G\) 成立都蕴含 \(a \leqslant 0\)。（为证明该条件充分，考虑形如

\[
x \mapsto \sum_ {i = 1} ^ {n} (f _ {i} (x) - f _ {i} (g _ {i} x h _ {i}))
\]

的函数所成的集合，并应用 Hahn--Banach 定理。）

\begin{enumerate}
\def\labelenumi{\alph{enumi})}
\setcounter{enumi}{1}
\item
  每个紧拓扑群都是可和的，每个阿贝尔拓扑群也都是可和的。
\item
  令 G 为非交换自由群。离散群 G 不是可和的。
\item
  对每个希尔伯特空间 E 及 G 在 E 中的每个连续有界表示，E 上都存在一个希尔伯特范数 q，使得 G 在赋予 q 的希尔伯特空间 E 中的表示是酉的。
\end{enumerate}

\begin{enumerate}
\def\labelenumi{\arabic{enumi})}
\setcounter{enumi}{4}
\tightlist
\item
  令 G 为拓扑群，H 为 G 的开子群。记 p 为典范投影 \(G \rightarrow G/H\)。
\end{enumerate}

\begin{enumerate}
\def\labelenumi{\alph{enumi})}
\item
  令 \(\sigma\colon\mathrm{G/H}\to\mathrm{G}\) 为满足 \(p\circ\sigma=\mathrm{Id}_{\mathrm{G/H}}\) 的映射。记 \(\beta\colon\mathrm{G}\times\mathrm{G/H}\to\mathrm{H}\) 为由 \(\beta(g,x)=\sigma(gx)^{-1}g\) 定义的映射。映射 \(\beta\) 是连续的。
\item
  令 \(\pi\) 为 H 在希尔伯特空间 E 中的有界表示。记 F 为希尔伯特空间 \(\ell^2 (\mathrm{G / H};\mathrm{E})\)。公式
\end{enumerate}

\[
\varrho (g) f (x) = \pi (\beta (g ^ {- 1}, x) ^ {- 1}) f (g ^ {- 1} x)
\]

定义了 G 在 F 中的连续有界表示。

\begin{enumerate}
\def\labelenumi{\alph{enumi})}
\setcounter{enumi}{2}
\tightlist
\item
  若 F 上存在一个希尔伯特范数 q，使得表示 \(\varrho\) 是 G 在 \((F,q)\) 中的酉表示，则 E 上存在一个希尔伯特范数 \(\widetilde{q}\)，使得 \(\pi\) 是 H 在 \((E,\widetilde{q})\) 中的酉表示。
\end{enumerate}

\begin{enumerate}
\def\labelenumi{\arabic{enumi})}
\setcounter{enumi}{5}
\tightlist
\item
  \begin{enumerate}
  \def\labelenumii{\alph{enumii})}
  \tightlist
  \item
    令 E 为复巴拿赫空间，令 \(u \in \mathcal{L}(\mathrm{E})\)。若存在 \(n \in N\)，使得 \(n \geqslant 1\) 且 \(u^{n} = 1_{E}\)，则 u 的谱包含于 n 次单位根群，并由 u 的特征值组成。b) 令 G 为交换拓扑群，且 G 的每个元素都有有限阶。G 在巴拿赫空间中的每个连续不可约表示维数均为 1，且取值于 \(U \subset C^{*}\)——模为 1 的复数所成的群。
  \end{enumerate}
\end{enumerate}

\begin{enumerate}
\def\labelenumi{\alph{enumi})}
\setcounter{enumi}{2}
\tightlist
\item
  令 G 为交换拓扑群，且 G 的有限阶元素集在 G 中稠密。G 在巴拿赫空间中的每个连续不可约表示维数均为 1，且取值于 \(U \subset C^{*}\)——模为 1 的复数所成的群。
\end{enumerate}

\begin{enumerate}
\def\labelenumi{\arabic{enumi})}
\setcounter{enumi}{6}
\tightlist
\item
  对每个整数 \(n \geqslant 1\)，记 \(s_{n}\) 为由下式定义的函数 \(Z \rightarrow R\)
\end{enumerate}

\[
s _ {n} (k) = \left\{ \begin{array}{l l} (1 - k) ^ {n (1 - k)} & \text{若 } k \leqslant - 1 \\ 1 & \text{若 } k = 0 \\ (k + 1) ^ {- (k + 1) / n} & \text{若 } k \geqslant 1. \end{array} \right.
\]

记 S 为函数 \(s_n\) 所成的集合。

\begin{enumerate}
\def\labelenumi{\alph{enumi})}
\item
  对每个整数 \(n \geqslant 1\) 及每个 \((k_{1}, k_{2}) \in \mathbf{Z}^{2}\)，有 \(s_{4n}(k_{1})s_{4n}(k_{2}) \geqslant s_{n}(k_{1} + k_{2})\)。
\item
  对于 \(f \in \mathbf{C}(\mathrm{T})\)，记 \((a_k(f))_{k \in \mathbf{Z}}\) 为 \(f\) 在 0 处的 Laurent 展开的系数。对于每个 \(s \in \mathrm{S}\) 和每个 \(f \in \mathbf{C}(\mathrm{T})\)，级数
\end{enumerate}

\[
\sum_ {k \in \mathbf {Z}} | a _ {k} (f) | s (k)
\]

收敛，并且集合

\[
\mathcal {V} (s, \varepsilon) = \{f \in \mathbf {C} (\mathrm{T}) | \sum_ {k \in \mathbf {Z}} | a _ {k} (f) | s (k) <   \varepsilon \}
\]

对 \(s \in S\) 和 \(\varepsilon > 0\) 定义的这些集合，构成 \(\mathbf{C}(T)\) 上某个局部凸拓扑的 0 的开邻域基本系。乘法 \(\mathbf{C}(T) \times \mathbf{C}(T) \to \mathbf{C}(T)\) 关于此拓扑是连续的。

\begin{enumerate}
\def\labelenumi{\alph{enumi})}
\setcounter{enumi}{2}
\tightlist
\item
  上一问题中定义的拓扑可度量化但不完备。
\end{enumerate}

\(d)\) 对于 \(f \in \mathbf{C}(\mathrm{T})^*\)，记 \(\varrho(f)\) 为映射 \(g \mapsto fg\)，从 \(\mathbf{C}(\mathrm{T})\) 到 \(\mathbf{C}(\mathrm{T})\)。映射 \(\varrho\) 是 \(\mathbf{C}(\mathrm{T})^*\) 在局部凸空间 \(\mathbf{C}(\mathrm{T})\) 中的连续表示。

\begin{enumerate}
\def\labelenumi{\alph{enumi})}
\setcounter{enumi}{4}
\tightlist
\item
  表示 \(\varrho\) 是不可约的。
\end{enumerate}

\exercisesection{}

在本节习题中，除非另有说明，G 表示赋有左哈尔测度的局部紧拓扑群，该测度记为 \(\mu\)。

\begin{enumerate}
\def\labelenumi{\arabic{enumi})}
\item
  令 \(\pi\) 为 G 在复希尔伯特空间 E 上的酉表示。表示 \(\pi\) 不可约，当且仅当 \(\mathcal{H}(\mathrm{G})\) 在 \(\mathcal{L}(\mathrm{E})\) 中的像关于逐点收敛拓扑是稠密的。（使用 IV, 第 326 页习题 24。）
\item
  令 H 为 G 的闭子群。令 \(\nu\) 为 G/H 上的拟不变测度（INT, VII, 第 56 页，第 2 节第 5 号定理 2），并令 \(\varrho\) 为与之关联的、取值于 \(\mathbf{R}_{+}^{*}\) 的 G 上函数。由下式定义的映射 \(\varrho(g)\)
\end{enumerate}

\[
\varrho (g) f (x) = \left(\frac {\varrho (g ^ {- 1} x)}{\varrho (x)}\right) ^ {1 / 2} f (g ^ {- 1} x)
\]

定义了 G 在 \(\mathcal{L}^{2}(\mathrm{G / H},\nu)\) 上的连续线性表示，并通过取商定义了 G 在 \(\mathrm{L}^2 (\mathrm{G / H},\nu)\) 上的酉表示，称为 \(\mathrm{G / H}\) 的拟正则表示。

\begin{enumerate}
\def\labelenumi{\arabic{enumi})}
\setcounter{enumi}{2}
\tightlist
\item
  令 G = R，令 H 为赋予离散拓扑并配备计数测度的群 R；该测度是哈尔测度。
\end{enumerate}

\begin{enumerate}
\def\labelenumi{\alph{enumi})}
\item
  对每个 \(t \in G\) 及每个 \(f \in L^{2}(H)\)，映射 \(\varrho(t)f: x \mapsto f(x + t)\) 属于 \(L^{2}(H)\)；映射 \(\varrho(t)\) 是 \(L^{2}(H)\) 上的酉映射。b) 对所有 \(f,g \in L^{2}(H)\)，由 \(t \mapsto \langle f | \varrho(t)g \rangle\) 定义的从 G 到 C 的映射是可测的。
\item
  令 \(f\) 和 \(g\) 属于 \(\mathrm{L}^2(\mathrm{H})\)。若存在 \(t \in \mathbf{R}\) 使得 \(\langle f | \varrho(t)g \rangle\) 非零，则映射 \(t \mapsto \langle f | \varrho(t)g \rangle\) 从 \(\mathbf{G}\) 到 \(\mathbf{C}\) 不连续。（特别地，E 可数生成这一假设不能从 INT, VIII, 第 171 页第 4 节第 7 号推论 2 中省略。）
\end{enumerate}

\(d)\) 存在 \(f \in \mathrm{L}^2(\mathrm{H})\)，使得 \(\mathrm{L}^2(\mathrm{H})\) 中由元素 \(\varrho(t)f\)（其中 \(t \in \mathbf{R}\)）生成的子空间在 \(\mathrm{L}^2(\mathrm{H})\) 中稠密。

\begin{enumerate}
\def\labelenumi{\arabic{enumi})}
\setcounter{enumi}{3}
\item
  令 \(p \geqslant 1\)。G 在 \(\mathrm{L}^{p}(\mathrm{G}, \mu)\) 上的双正则表示的核，是 G 的中心在同态 \(g \mapsto (g, g)\) 下的像。
\item
  令 G 为局部紧单模群。
\end{enumerate}

\begin{enumerate}
\def\labelenumi{\alph{enumi})}
\item
  若存在 G 的有限维平方可积表示，则 G 是紧的。
\item
  若存在 G 的平方可积表示，则 G 的中心是紧的。
\end{enumerate}

\begin{enumerate}
\def\labelenumi{\arabic{enumi})}
\setcounter{enumi}{5}
\tightlist
\item
  假设群 G 是单模的。记 \(\widehat{G}_{d} \subset \widehat{G}\) 为 G 的不可约酉平方可积表示的等价类集。对于 \(\pi \in \widehat{G}_{d}\)，记 \(p_{\pi}\) 为 \(L^{2}(G)\) 上的正交投影，其像是 G 的右正则表示的 \(\pi\)-同型分量。
\end{enumerate}

\(a)\) 令 \(f \in \mathcal{L}^2(\mathrm{G})\) 且 \(\pi \in \widehat{\mathrm{G}}_d\)。记 E 为 \(\pi\) 的表示空间。映射 \(g \mapsto f(g)\pi(g)\) 属于 \(\mathrm{L}^1(\mathrm{G};\mathcal{L}(\mathrm{E}))\)；自同态

\[
v = \int_ {\mathrm{G}} f (g) \pi (g) d \mu (g)
\]

是希尔伯特–施密特自同态，记为 \(v = \pi(f)\)。

\begin{enumerate}
\def\labelenumi{\alph{enumi})}
\setcounter{enumi}{1}
\item
  令 \(f \in \mathcal{L}^2(\mathrm{G})\)。有 \(\| \pi(\overline{f}) \|_2 = \| p_\pi(f) \|\)。
\item
  令 \(L_{d}^{2}(G)\) 为 \(L^{2}(G)\) 的闭子空间，它由 \(\pi\)-同型分量生成，这些分量属于 \(L^{2}(G)\)，且 \(\pi \in \widehat{G}_{d}\)。它是各正交投影 \(p_{\pi}\) 的像（对 \(\pi \in \widehat{G}_{d}\)）的希尔伯特直和。对于每个 \(f \in L_{d}^{2}(G)\)，有
\end{enumerate}

\[
\| f \| ^ {2} = \sum_ {\pi \in \widehat {\mathrm{G}} _ {d}} c (\pi) \| \pi (f) \| _ {2} ^ {2}
\]

其中 \(c(\pi)\) 是 \(\pi\) 相对于 \(\{e\}\) 的形式次数。

\begin{enumerate}
\def\labelenumi{\arabic{enumi})}
\setcounter{enumi}{6}
\tightlist
\item
  令 \(n \in N\)。确定 \(R^{n}\) 的酉自守表示（即对每个离散子群 \(\Gamma\)，它们都是 \(\Gamma\)-自守的 \(R^{n}\) 的表示）。
\end{enumerate}

8) 假设 G 是局部紧单模群。令 Γ 为 G 的离散子群，使商空间 X = Γ\textbackslash G 是紧的。令 μ 为 G 上的哈尔测度，β 为 Γ 上的计数测度；记 \(\widetilde{\mu} = \mu/\beta\)。

令 \(\varrho\) 为 V, 第 404 页例 3 中定义的 G 在 \(\mathrm{L}^{2}(\mathrm{X},\widetilde{\mu})\) 上的酉表示。

对于 \(\gamma\in\Gamma\)，记 \(\Gamma_{\gamma}\)（分别为 \(\mathbf{G}_{\gamma}\)）为 \(\gamma\) 在 \(\Gamma\)（分别在 \(\mathbf{G}\)）中的中心化子。\(a)\) 群 \(\mathbf{G}_{\gamma}\) 是单模的。

\begin{enumerate}
\def\labelenumi{\alph{enumi})}
\setcounter{enumi}{1}
\item
  令 \(\varphi\in\mathcal{K}(\mathrm{G})\)。对于每个 \(\Gamma\)-自守酉表示 \(\pi\)，自同态 \(\pi(\varphi)\) 都有有限迹。
\item
  有
\end{enumerate}

\[
\operatorname{Tr} (\varrho (\varphi)) = \sum_ {\pi \in \widehat {\mathrm{G}}} m _ {\pi} (\varrho) \operatorname{Tr} (\pi (\varphi))
\]

其中 \(m_{\pi}(\varrho)\) 是 \(\pi\) 在 \(\varrho\) 中的重数（“迹公式的谱侧”）。

\begin{enumerate}
\def\labelenumi{\alph{enumi})}
\setcounter{enumi}{3}
\tightlist
\item
  令 \(\Gamma^{\sharp}\) 为 \(\Gamma\) 的共轭类集。对于 \(\gamma \in \Gamma^{\sharp}\)，记 \(\beta_{\gamma}\) 为 \(\Gamma_{\gamma}\) 上的计数测度。有
\end{enumerate}

\[
\mathrm{Tr} (\varrho (\varphi)) = \sum_ {\gamma \in \Gamma^ {\sharp}} \int_ {\Gamma_ {\gamma} \backslash G} \varphi (x ^ {- 1} \gamma x) d (\mu / \beta_ {\gamma}) (x).
\]

\begin{enumerate}
\def\labelenumi{\alph{enumi})}
\setcounter{enumi}{4}
\tightlist
\item
  固定 \(\mu_{\gamma}\) 在 \(G_{\gamma}\) 上的哈尔测度，对所有 \(\gamma \in \Gamma\)。有
\end{enumerate}

\[
\operatorname{Tr} (\varrho (\varphi)) = \sum_ {\gamma \in \Gamma^ {\sharp}} \int_ {G _ {\gamma} \backslash G} \int_ {\Gamma_ {\gamma} \backslash G _ {\gamma}} \varphi (x ^ {- 1} y ^ {- 1} \gamma y x) d (\mu_ {\gamma} / \beta_ {\gamma}) (y) d (\mu / \mu_ {\gamma}) (x).
\]

\(f)\) 对每个 \(\gamma \in \Gamma\)，测度 \(\mu_{\gamma} / \beta_{\gamma}\) 在 \(\Gamma_{\gamma} \backslash G_{\gamma}\) 上是有限的。\(g)\) 有

\[
\mathrm{Tr} (\varrho (\varphi)) = \sum_ {\gamma \in \Gamma^ {\sharp}} (\mu_ {\gamma} / \beta_ {\gamma}) (\Gamma_ {\gamma} \backslash \mathrm{G} _ {\gamma}) \int_ {\mathrm{G} _ {\gamma} \backslash \mathrm{G}} \varphi (x ^ {- 1} \gamma x) d (\mu / \mu_ {\gamma}) (x)
\]

（“迹公式的几何侧”；将此表达式与 c) 中所得表达式相等联系起来的公式，称为 Γ-自守表示的“迹公式”。）

\begin{enumerate}
\def\labelenumi{\alph{enumi})}
\setcounter{enumi}{7}
\tightlist
\item
  置 G = R 且 \(\Gamma = Z\)。验证此时迹公式等价于泊松公式（参见 II, 第 239 页）。
\end{enumerate}

\begin{enumerate}
\def\labelenumi{\arabic{enumi})}
\setcounter{enumi}{8}
\item
  假设 G 是无紧因子的完美实半单李群。G 没有非平凡的有限维酉表示。
\item
  \begin{enumerate}
  \def\labelenumii{\alph{enumii})}
  \tightlist
  \item
    G 在 \(L^{2}(G)\) 上的左正则表示的每个矩阵系数都属于 \(\mathcal{C}_{0}(G)\)。
  \end{enumerate}
\end{enumerate}

\begin{enumerate}
\def\labelenumi{\alph{enumi})}
\setcounter{enumi}{1}
\tightlist
\item
  若 G 非紧，则 G 在 L²(G) 上的左正则表示不含非零有限维子表示。
\end{enumerate}

\begin{enumerate}
\def\labelenumi{\arabic{enumi})}
\setcounter{enumi}{10}
\tightlist
\item
  令 E 为希尔伯特空间。
\end{enumerate}

\begin{enumerate}
\def\labelenumi{\alph{enumi})}
\item
  令 \(\varrho\) 为群 \(\mathbf{R}\) 在 E 中的酉表示。令 D 为 E 的稠密子空间，在 \(\varrho\) 下稳定，并且对每个 \(x\in\mathrm{D}\)，映射 \(\psi_{x}:t\mapsto\varrho(t)x\) 在 R 上可微。令 u 为定义域为 D 且满足 \(u(x) = i^{-1}\psi_x'(0)\)（对所有 \(x \in D\)）的部分算子。则 u 本质自伴，且 \(\overline{u}\) 是 \(\varrho\) 的无穷小生成元。
\item
  令 u 为 E 上的自伴部分算子，D 为 \(\text{dom}(u)\) 的稠密子空间。假设对每个 \(t \in R\)，有 \(\exp(itu)(D) \subset D\)。则空间 D 在 \(E_{u}\) 中稠密。
\item
  令 u 为 E 上的自伴部分算子。令 F 为希尔伯特空间，v 为 F 上的自伴部分算子。部分算子 \(u \otimes 1_{F} + 1_{E} \otimes E\) 在 \(E \otimes F\) 上定义了 \(E \widehat{\otimes}_{2} F\) 上的一个本质自伴部分算子。
\end{enumerate}

\begin{enumerate}
\def\labelenumi{\arabic{enumi})}
\setcounter{enumi}{11}
\tightlist
\item
  令 G 为实李群，令 \(\varrho\) 为 G 在希尔伯特空间 E 中的酉表示。若元素 \(x \in \mathrm{E}\) 为 \(\mathrm{C}^{\infty}\) 类，亦即从 G 到 E 的映射 \(g \mapsto \varrho(g)x\) 为 \(\mathrm{C}^{\infty}\) 类，则称其为类向量。a) 令 \(\pi\) 为 G 在 F 中的酉表示，且 \(u: \mathrm{E} \to \mathrm{F}\) 为 G-态射。对于每个 \(x \in \mathrm{E}\) 的 \(\mathrm{C}^{\infty}\) 类向量，向量 \(u(x)\) 在 F 中也是 \(\mathrm{C}^{\infty}\) 类的。
\end{enumerate}

记 \(\mathrm{C}^{\infty}(\varrho)\) 为 E 中 \(C^{\infty}\) 类向量的空间。由问题 a)，每个从 \(\varrho\) 到 \(\pi\) 的 G-态射通过取子空间定义了从 \(\mathrm{C}^{\infty}(\varrho)\) 到 \(\mathrm{C}^{\infty}(\pi)\) 的线性映射。

\begin{enumerate}
\def\labelenumi{\alph{enumi})}
\setcounter{enumi}{1}
\tightlist
\item
  空间 \(\mathrm{C}^{\infty}(\varrho)\) 在 E 中稠密。
\end{enumerate}

\begin{enumerate}
\def\labelenumi{\arabic{enumi})}
\setcounter{enumi}{12}
\tightlist
\item
  令 \(\pi\) 为 G 在复希尔伯特空间 E 中的不可约酉表示。令 \(\varrho\) 为 G 在复希尔伯特空间 F 中的酉表示。假设存在一个从 E 到 F 的闭稠定部分算子 \(u\)，其定义域是 E 的 G-稳定子空间，并且满足 \(u \circ \pi(g) = \varrho(g) \circ u\) 对所有 \(g \in \mathbf{G}\) 成立。
\end{enumerate}

部分算子 \(u\) 属于 \(\mathcal{L}(\mathrm{E};\mathrm{F})\)，并且存在 \(\lambda > 0\)，使得 \(\lambda u\) 是等距的。

\begin{enumerate}
\def\labelenumi{\arabic{enumi})}
\setcounter{enumi}{13}
\tightlist
\item
  令 Z 为 G 的中心 C 的子群，且 C/Z 是紧的；令 \(\nu\) 为 G/Z 上的哈尔测度。令 \(\pi\) 为 G 在复希尔伯特空间 E 中的不可约酉表示。假设 \(\pi\) 模中心平方可积。令 \(u \in \mathcal{L}(\mathrm{E})\) 为核型自同态（IV, 第 169 页定义 4）。
\end{enumerate}

\begin{enumerate}
\def\labelenumi{\alph{enumi})}
\tightlist
\item
  对 E 中的每个 \(x\) 和 \(y\)，函数
\end{enumerate}

\[
g \mapsto \langle x | \pi (g) u \pi (g ^ {- 1}) y \rangle
\]

属于 \(\mathcal{F}(\mathrm{G / Z})\)。

\begin{enumerate}
\def\labelenumi{\alph{enumi})}
\setcounter{enumi}{1}
\tightlist
\item
  对 E 中的每个 \(x\) 和 \(y\)，有
\end{enumerate}

\[
\int_ {\mathrm{G} / \mathrm{Z}} \langle x \mid \pi (g) u \pi (g ^ {- 1}) y \rangle d \nu (g) = \frac {1}{c _ {\mathrm{Z}} (\pi)} \langle x \mid y \rangle \operatorname{Tr} (u).
\]

\begin{enumerate}
\def\labelenumi{\arabic{enumi})}
\setcounter{enumi}{14}
\tightlist
\item
  令 Z 为 G 的中心 C 的子群，且 C/Z 是紧的；令 \(\nu\) 为 G/Z 上的哈尔测度。令 \(\pi\) 为 G 在复希尔伯特空间 E 中的不可约酉表示。记 \(\chi\) 为 \(\pi\) 的中心特征标在 Z 上的限制。
\end{enumerate}

记 \(f_{x,y}\) 为矩阵系数 \(g \mapsto \langle x | \pi(g)y \rangle\)，其中 \((x, y) \in \mathrm{E}^{2}\)。令 \(p \geqslant 1\) 为实数，令 \(q\) 为 p 的共轭指数。

\begin{enumerate}
\def\labelenumi{\alph{enumi})}
\item
  令 \(y \in \mathrm{E}\)，使得矩阵系数 \(f_{x,y}\) 属于 \(\mathcal{L}_{\chi}^{p}(\mathrm{G})\)。映射 \(x \mapsto f_{x,y}\) 是从 \(\pi\) 到 \(\boldsymbol{\delta}_{\mathrm{G},\chi}^{(p)}\) 的连续 G-态射。
\item
  令 \(\varrho\) 为 G 在复希尔伯特空间 F 中的不可约酉表示，其中心特征标限制所得的 \(\chi\) 定义在 \(Z\) 上。假设 \(p > 1\)，且 \(\pi\) 的矩阵系数属于 \(\mathcal{L}_{\chi}^{p}(\mathrm{G})\)，\(\varrho\) 的矩阵系数属于 \(\mathcal{L}_{\chi}^{q}(\mathrm{G})\)。令 \(x_0 \in \mathrm{E}\) 且 \(y_0 \in \mathrm{F}\)。对于每个 \(x \in \mathrm{E}\)，存在从 F 到 C 的连续线性泛函 \(\ell_x\)，使得
\end{enumerate}

\[
\ell_ {x} (y) = \int_ {\mathrm{G/Z}} \overline {{\langle x _ {0} | \pi (g) x \rangle}} \langle y _ {0} | \varrho (g) y \rangle d \nu (g)
\]

对所有 \(y \in \mathbf{F}\) 都成立。

\begin{enumerate}
\def\labelenumi{\alph{enumi})}
\setcounter{enumi}{2}
\item
  存在从 E 到 F 的连续线性映射 u，使得有 \(\ell_{x}(y)=\langle u(x)\mid y\rangle\)，对所有 \((x,y)\in\mathrm{E}\times\mathrm{F}\) 成立。
\item
  有
\end{enumerate}

\[
\int_ {\mathrm{G} / \mathrm{Z}} \overline {{\langle x | \pi (g) x ^ {\prime} \rangle}} \langle y | \varrho (g) y ^ {\prime} \rangle d \nu (g) = 0
\]

对 \((x,x^{\prime},y,y^{\prime})\in \mathrm{E}^{2}\times \mathrm{F}^{2}\) 都成立。

\begin{enumerate}
\def\labelenumi{\alph{enumi})}
\setcounter{enumi}{4}
\tightlist
\item
  假设 p = 1。\(\pi\) 的矩阵系数与 \(\varrho\) 的矩阵系数正交。
\end{enumerate}

\begin{enumerate}
\def\labelenumi{\arabic{enumi})}
\setcounter{enumi}{15}
\tightlist
\item
  令 \(\mathbf{G} = \mathbf{S}\mathbf{L}(2, \mathbf{R})\)，令 H 为 C 中由满足 \(\mathcal{I}(z) > 0\) 的复数 z 组成的开子集。令 \(\mu\) 为 C 上的勒贝格测度。对于 G 中的每个元素
\end{enumerate}

\[
g = \left( \begin{array}{c c} a & b \\ c & d \end{array} \right)
\]

及每个 \(z \in \mathbf{H}\)，置 \(j(g, z) = cz + d\)。

\begin{enumerate}
\def\labelenumi{\alph{enumi})}
\tightlist
\item
  群 G 通过下式在 H 上连续地左作用
\end{enumerate}

\[
\left( \begin{array}{c c} a & b \\ c & d \end{array} \right) \cdot z = \frac {a z + b}{c z + d}.
\]

测度 \(\nu = y^{-2} \cdot \mu\) 是 G-不变的。

\begin{enumerate}
\def\labelenumi{\alph{enumi})}
\setcounter{enumi}{1}
\tightlist
\item
  令 n 为满足 \(n \geqslant 2\) 的整数。记 \(E_{n}\) 为全纯函数 \(f: H \to C\) 所成的向量空间，其中函数 \(z \mapsto |f(z)|^{2} \mathcal{I}(z)^{n}\) 属于 \(\mathcal{L}^{2}(\mathbf{H}, \nu)\)。空间 \(E_{n}\) 在赋予半双线性形式
\end{enumerate}

\[
\langle f \mid g \rangle = \int_ {\mathbf {H}} \overline {{f (z)}} g (z) \mathcal {I} (z) ^ {n} d \nu (z)
\]

后，是可数型且无限维的希尔伯特空间。

\begin{enumerate}
\def\labelenumi{\alph{enumi})}
\setcounter{enumi}{2}
\tightlist
\item
  对每个 \(g \in G\) 及每个 \(f \in E_{n}\)，置
\end{enumerate}

\[
\pi_ {n} (g) f (z) = j (g ^ {- 1}, z) ^ {- n} f (g ^ {- 1} \cdot z).
\]

映射 \(\pi_n\) 定义了 G 在 \(\mathrm{E}_n\) 中的酉表示。\(d)\) 由函数 \(f_n\) 定义为 \(f_n(z) = (z + i)^{-n}\)，它属于 \(\mathrm{E}_n\)；它满足 \(\pi_n(k)f_n = \chi_{-n}(k)f_n\) 对所有 \(k \in \mathbf{SO}(\mathbf{R}^2)\) 成立，其中 \(\chi_n\) 是满足下式的特征标

\[
\chi_ {n} \Big (\left( \begin{array}{c c} \cos (\theta) & \sin (\theta) \\ - \sin (\theta) & \cos (\theta) \end{array} \right) \Big) = e ^ {i n \theta}
\]

对所有 \(\theta\in\mathbf{R}\) 都成立。

\begin{enumerate}
\def\labelenumi{\alph{enumi})}
\setcounter{enumi}{4}
\item
  令 F 为 \(E_{n}\) 的子表示。\(E_{n}\) 的正交投影，其像是 \(\chi_{-n}\)-同型分量，属于将 \(\pi_{n}\) 限制到 \(\mathbf{SO}(\mathbf{R}^{2})\) 后的表示，由 \(f \mapsto -(2i)^{n}f(i)f_{n}\) 给出。（使用 V, 第 465 页推论 2 以及柯西积分公式，FVC，即将出版。）
\item
  表示 \(\pi_n\) 是不可约的。
\end{enumerate}

\(g)\) 矩阵系数 \(g \mapsto \langle f_n | \pi_n(g)f_n \rangle\) 非零且属于 \(\mathcal{L}^2(\mathrm{G})\)。

\begin{enumerate}
\def\labelenumi{\alph{enumi})}
\setcounter{enumi}{7}
\tightlist
\item
  表示 \(\pi_{n}\) 是平方可积的，并且存在常数 c > 0，使得 \(\pi_{n}\) 的形式次数为 \(c(n - 1)\)。
\end{enumerate}

\begin{enumerate}
\def\labelenumi{\arabic{enumi})}
\setcounter{enumi}{16}
\item
  令 G 为连通实李群。G 存在忠实有限维酉表示，当且仅当群 G 局部同构于紧群。
\item
  令 X 为赋有有界正测度 \(\mu\) 的局部紧空间。对于每个巴拿赫空间 F，记 S(X, \(\mu\) ; F) 为从 X 到 F 的 \(\mu\)-可测函数按 \(\mu\)-几乎处处相等关系所得的等价类空间，并赋予依测度收敛拓扑（INT, IV, 第 194 页，第 5 节第 11 号）。记 A 为 \(\mathcal{L}(\mathrm{L}_{\mathbf{C}}^{2}(\mathrm{X},\mu))\) 的交换子代数，它由函数 \(f \in \mathcal{L}_{\mathbf{C}}^{\infty}(\mathrm{X},\mu)\) 所对应的乘法自同态组成。
\end{enumerate}

记 \(e\colon \mathbf{R}\to \mathbf{C}\) 为同态 \(t\mapsto \exp (2i\pi t)\)。

令 E 为拓扑向量空间。

\begin{enumerate}
\def\labelenumi{\alph{enumi})}
\tightlist
\item
  令 \(f \in \mathrm{S}(\mathrm{X}, \mu; \mathbf{R})\)，使得 \(|f(x)| \geqslant 1\) 对 \(\mu\)-几乎所有 \(x \in \mathrm{X}\) 都成立。存在实数 \(t \in ]0, 1[\)，使得
\end{enumerate}

\[
\mu (\{x \in \mathrm{X} \mid t f (x) \in [ 1 / 4, 3 / 4 ] + \mathbf {Z} \}) \geqslant \frac {2}{5} \mu (\mathrm{X}).
\]

\begin{enumerate}
\def\labelenumi{\alph{enumi})}
\setcounter{enumi}{1}
\item
  令 u 为从 E 到 \(\mathrm{S}(\mathrm{X},\mu;\mathbf{R})\) 的线性映射。对于每个 \(x\in E\)，记 \(\varrho_{u}(x)\) 为 A 中由函数 \(e\circ u(x)\) 的乘法自同态给出的元素，该函数属于 \(\mathcal{L}_{\mathbf{C}}^{\infty}(\mathrm{X},\mu)\)。映射 \(\varrho_{u}\) 是从 G 到 \(\mathrm{L}_{\mathbf{C}}^{2}(\mathrm{X},\mu)\) 的连续表示，当且仅当映射 u 连续。（若 u 不连续，使用上一问题验证 \(\varrho\) 不连续。）
\item
  令 \(\varrho\) 为 E 在 \(\mathrm{L}^2 (\mathrm{X},\mu)\) 中的酉表示，且其像包含于 \(\mathcal{A}\)。存在唯一的连续线性映射 \(u\colon \mathrm{E}\to \mathrm{S}(\mathrm{X},\mu ;\mathbf{R})\)，使得 \(\varrho = \varrho_{u}\)。（先考虑 \(\mathrm{E} = \mathbf{R}\) 的情形，然后应用斯通定理。）
\end{enumerate}

\begin{enumerate}
\def\labelenumi{\arabic{enumi})}
\setcounter{enumi}{18}
\item
  令 G 为有限群，H 为 G 的子群。对于 H 在有限维复希尔伯特空间 E 中的每个酉表示 \(\pi\)，证明诱导酉表示 \(\operatorname{Ind}_{\mathrm{H}}^{\mathrm{G}}(\pi)\) 同构于 A, VIII, 第 392 页第 21 节第 3 号中定义的诱导表示。
\item
  令 G 为无穷远处可数的局部紧群，令 H 为 G 的闭子群。令 \(\pi\) 为 H 在希尔伯特空间 E 中的酉表示，令 \(\varrho\) 为由 \(\pi\) 诱导的 G 的酉表示。记 F 为 \(\varrho\) 的表示空间。
\end{enumerate}

在 G 上固定哈尔测度 \(\mu\)，在 H 上固定哈尔测度 \(\beta\)，并记 \(\kappa\colon\mathbf{G}\to\mathbf{R}_{+}^{*}\) 为满足 \(\kappa(xh)=\Delta_{\mathrm{H}}(h)\Delta_{\mathrm{G}}(h)^{-1}\kappa(x)\) 对 \((x,h)\in\mathbf{G}\times\mathrm{H}\) 成立，且满足 \(\kappa(e)=1\) 的连续函数。

\begin{enumerate}
\def\labelenumi{\alph{enumi})}
\item
  若 G/H 具有有界 G-不变测度，且存在非零的 H-不变向量 \(x \in E\)，则存在非零的 G-不变向量 \(y \in F\)。
\item
  反之，设 \(y \in F\) 是 \(\varrho\) 的非零 G-不变向量。证明存在非零函数 \(f \in L_{\overline{\pi}}^{2}\)，使得
\end{enumerate}

\[
f (x) = \left(\frac {\kappa (g ^ {- 1} x)}{\kappa (x)}\right) ^ {1 / 2} f (g ^ {- 1} x)
\]

对所有 \((g, x) \in \mathrm{G}^{2}\) 都成立。推出 \(\kappa = 1\)，从而 G/H 上的测度 \(\mu/\beta\) 有定义且有界，最后证明 \(f(e) \in \mathrm{E}\) 是非零 H-不变向量。